%% Supplementary material for "It Is Not the Signature: An Empirical Study of a
%% Hidden Key-Parse Cost in JSON Web Token Validation" (Journal of Systems and
%% Software). Compile main.tex first: cross-references to the main paper are
%% resolved from main.aux through the xr-hyper package.
\documentclass[preprint,12pt,authoryear]{elsarticle}
\usepackage{lmodern}
\usepackage{amsmath}
\usepackage{amssymb}
\usepackage{graphicx}
\usepackage{booktabs}
\usepackage{array}
\usepackage{xurl}
\usepackage{xcolor}
\usepackage{xr-hyper}
\usepackage[hidelinks]{hyperref}
\externaldocument{main}
\IfFileExists{microtype.sty}{\usepackage{microtype}}{}
% keypath_macros.tex -- GENERATED. Do not edit.
% Regenerate: .venv/bin/python -m analysis.make_keypath_macros
% Every value below is read from a file under data/runs/ or survey/.
\providecommand{\PLACEHOLDER}[1]{\textcolor{red}{\textbf{[MISSING: #1]}}}

\newcommand{\AbCpuDelta}{76.92}% data/runs/INSITU-AB-authmode-jwt/openloop_ramp.csv
\newcommand{\AbCpuJwt}{146.38}% data/runs/INSITU-AB-authmode-jwt/openloop_ramp.csv
\newcommand{\AbCpuNone}{69.45}% data/runs/INSITU-AB-authmode-none/openloop_ramp.csv
\newcommand{\AbCpuPerCallDelta}{384.6}% data/runs/INSITU-AB-authmode-jwt/openloop_ramp.csv
\newcommand{\AbCpuPerCallJwt}{731.9}% data/runs/INSITU-AB-authmode-jwt/openloop_ramp.csv
\newcommand{\AbCpuPerCallNone}{347.3}% data/runs/INSITU-AB-authmode-none/openloop_ramp.csv
\newcommand{\AbCpuSampleMaxJwt}{183.299}% data/runs/INSITU-AB-authmode-jwt/pod_cpu_samples.csv
\newcommand{\AbCpuSampleMaxNone}{100.577}% data/runs/INSITU-AB-authmode-none/pod_cpu_samples.csv
\newcommand{\AbCpuSampleMinJwt}{114.963}% data/runs/INSITU-AB-authmode-jwt/pod_cpu_samples.csv
\newcommand{\AbCpuSampleMinNone}{53.088}% data/runs/INSITU-AB-authmode-none/pod_cpu_samples.csv
\newcommand{\AbCpuSamplesJwt}{127}% data/runs/INSITU-AB-authmode-jwt/pod_cpu_samples.csv
\newcommand{\AbCpuSamplesNone}{127}% data/runs/INSITU-AB-authmode-none/pod_cpu_samples.csv
\newcommand{\AbCpuVsPredictedPct}{0.8}% data/runs/INSITU-AB-authmode-jwt/openloop_ramp.csv
\newcommand{\AbEventLoopLagJwt}{3.159}% data/runs/INSITU-AB-authmode-jwt/openloop_ramp.csv
\newcommand{\AbEventLoopLagNone}{3.137}% data/runs/INSITU-AB-authmode-none/openloop_ramp.csv
\newcommand{\AbGapSeconds}{78}% data/runs/INSITU-AB-authmode-jwt/run_metadata.json
\newcommand{\AbLatencyDeltaUs}{353.0}% data/runs/INSITU-AB-authmode-jwt/openloop_ramp.csv
\newcommand{\AbLatencyJwt}{2.0332}% data/runs/INSITU-AB-authmode-jwt/openloop_ramp.csv
\newcommand{\AbLatencyNone}{1.6802}% data/runs/INSITU-AB-authmode-none/openloop_ramp.csv
\newcommand{\AbLatencyPNineNineJwt}{8.0890}% data/runs/INSITU-AB-authmode-jwt/openloop_ramp.csv
\newcommand{\AbLatencyPNineNineNone}{6.3480}% data/runs/INSITU-AB-authmode-none/openloop_ramp.csv
\newcommand{\AbLoadFifteenJwt}{3.49}% data/runs/INSITU-AB-authmode-jwt/run_metadata.json
\newcommand{\AbLoadFifteenNone}{3.18}% data/runs/INSITU-AB-authmode-none/run_metadata.json
\newcommand{\AbLoadFiveJwt}{4.10}% data/runs/INSITU-AB-authmode-jwt/run_metadata.json
\newcommand{\AbLoadFiveNone}{3.50}% data/runs/INSITU-AB-authmode-none/run_metadata.json
\newcommand{\AbLoadOneJwt}{5.11}% data/runs/INSITU-AB-authmode-jwt/run_metadata.json
\newcommand{\AbLoadOneNone}{3.53}% data/runs/INSITU-AB-authmode-none/run_metadata.json
\newcommand{\AbRateJwt}{200.006}% data/runs/INSITU-AB-authmode-jwt/openloop_ramp.csv
\newcommand{\AbRateNone}{200.006}% data/runs/INSITU-AB-authmode-none/openloop_ramp.csv
\newcommand{\AbRequestsJwt}{36001}% data/runs/INSITU-AB-authmode-jwt/k6_summary.json
\newcommand{\AbRequestsNone}{36001}% data/runs/INSITU-AB-authmode-none/k6_summary.json
\newcommand{\AbSidecarCpuJwt}{49.13}% data/runs/INSITU-AB-authmode-jwt/openloop_ramp.csv
\newcommand{\AbSidecarCpuNone}{53.40}% data/runs/INSITU-AB-authmode-none/openloop_ramp.csv
\newcommand{\AbTopProcPctJwt}{84.3}% data/runs/INSITU-AB-authmode-jwt/run_metadata.json
\newcommand{\AbTopProcPctNone}{135.1}% data/runs/INSITU-AB-authmode-none/run_metadata.json
\newcommand{\AbWindowSecondsJwt}{180}% data/runs/INSITU-AB-authmode-jwt/k6_summary.json
\newcommand{\AbWindowSecondsNone}{180}% data/runs/INSITU-AB-authmode-none/k6_summary.json
\newcommand{\AdjudicatedFalsePositive}{5}% survey/data/hand_adjudication.csv
\newcommand{\AdjudicatedGenuine}{7}% survey/data/hand_adjudication.csv
\newcommand{\AdjudicatedGenuinePct}{5.36}% survey/data/hand_adjudication.csv
\newcommand{\AdjudicatedGenuineProjects}{6}% survey/data/hand_adjudication.csv
\newcommand{\AdjudicatedTotal}{13}% survey/data/hand_adjudication.csv
\newcommand{\AdjudicatedUnresolved}{1}% survey/data/hand_adjudication.csv
\newcommand{\ContainerVersusHostRatio}{0.92}% data/runs/2026-09-17T08-57-13Z-keypath-runtime-matrix/keypath_stats.json
\newcommand{\CounterFilesFetched}{191}% survey/data/counterexamples.json
\newcommand{\CounterReposWithVerify}{112}% survey/data/counterexamples.json
\newcommand{\CounterWithVerify}{120}% survey/data/counterexamples.json
\newcommand{\DecompExplainedPct}{77.9}% data/runs/2026-09-17T08-42-17Z-keypath-mechanism/keypath_stats.json
\newcommand{\DecompObserved}{23.92}% data/runs/2026-09-17T08-42-17Z-keypath-mechanism/keypath_stats.json
\newcommand{\DecompPredicted}{18.62}% data/runs/2026-09-17T08-42-17Z-keypath-mechanism/keypath_stats.json
\newcommand{\DecompResidual}{5.292}% data/runs/2026-09-17T08-42-17Z-keypath-mechanism/keypath_stats.json
\newcommand{\DecompResidualCiHi}{5.501}% data/runs/2026-09-17T08-42-17Z-keypath-mechanism/keypath_stats.json
\newcommand{\DecompResidualCiLo}{4.916}% data/runs/2026-09-17T08-42-17Z-keypath-mechanism/keypath_stats.json
\newcommand{\HostCallsPerInvocation}{40000}% data/runs/2026-09-17T08-42-17Z-keypath-mechanism/keypath_stats.json
\newcommand{\HostCreateSecretKey}{0.583}% data/runs/2026-09-17T08-42-17Z-keypath-mechanism/keypath_stats.json
\newcommand{\HostCreateSecretKeyCiHi}{0.583}% data/runs/2026-09-17T08-42-17Z-keypath-mechanism/keypath_stats.json
\newcommand{\HostCreateSecretKeyCiLo}{0.542}% data/runs/2026-09-17T08-42-17Z-keypath-mechanism/keypath_stats.json
\newcommand{\HostCvMax}{3.5}% data/runs/2026-09-17T08-42-17Z-keypath-mechanism/keypath_stats.json
\newcommand{\HostDecodeOnly}{1.292}% data/runs/2026-09-17T08-42-17Z-keypath-mechanism/keypath_stats.json
\newcommand{\HostDecodeOnlyCiHi}{1.292}% data/runs/2026-09-17T08-42-17Z-keypath-mechanism/keypath_stats.json
\newcommand{\HostDecodeOnlyCiLo}{1.250}% data/runs/2026-09-17T08-42-17Z-keypath-mechanism/keypath_stats.json
\newcommand{\HostHmacKeyObject}{2.166}% data/runs/2026-09-17T08-42-17Z-keypath-mechanism/keypath_stats.json
\newcommand{\HostHmacKeyObjectCiHi}{2.167}% data/runs/2026-09-17T08-42-17Z-keypath-mechanism/keypath_stats.json
\newcommand{\HostHmacKeyObjectCiLo}{2.125}% data/runs/2026-09-17T08-42-17Z-keypath-mechanism/keypath_stats.json
\newcommand{\HostHmacString}{2.208}% data/runs/2026-09-17T08-42-17Z-keypath-mechanism/keypath_stats.json
\newcommand{\HostHmacStringCiHi}{2.250}% data/runs/2026-09-17T08-42-17Z-keypath-mechanism/keypath_stats.json
\newcommand{\HostHmacStringCiLo}{2.208}% data/runs/2026-09-17T08-42-17Z-keypath-mechanism/keypath_stats.json
\newcommand{\HostInvocations}{15}% data/runs/2026-09-17T08-42-17Z-keypath-mechanism/keypath_stats.json
\newcommand{\HostJwtPreparsed}{6.042}% data/runs/2026-09-17T08-42-17Z-keypath-mechanism/keypath_stats.json
\newcommand{\HostJwtPreparsedCiHi}{6.125}% data/runs/2026-09-17T08-42-17Z-keypath-mechanism/keypath_stats.json
\newcommand{\HostJwtPreparsedCiLo}{5.958}% data/runs/2026-09-17T08-42-17Z-keypath-mechanism/keypath_stats.json
\newcommand{\HostJwtRsPemString}{52.58}% data/runs/2026-09-17T08-42-17Z-keypath-mechanism/keypath_stats.json
\newcommand{\HostJwtRsPemStringCiHi}{52.71}% data/runs/2026-09-17T08-42-17Z-keypath-mechanism/keypath_stats.json
\newcommand{\HostJwtRsPemStringCiLo}{52.42}% data/runs/2026-09-17T08-42-17Z-keypath-mechanism/keypath_stats.json
\newcommand{\HostJwtRsPreparsed}{25.33}% data/runs/2026-09-17T08-42-17Z-keypath-mechanism/keypath_stats.json
\newcommand{\HostJwtRsPreparsedCiHi}{25.50}% data/runs/2026-09-17T08-42-17Z-keypath-mechanism/keypath_stats.json
\newcommand{\HostJwtRsPreparsedCiLo}{25.29}% data/runs/2026-09-17T08-42-17Z-keypath-mechanism/keypath_stats.json
\newcommand{\HostJwtString}{29.92}% data/runs/2026-09-17T08-42-17Z-keypath-mechanism/keypath_stats.json
\newcommand{\HostJwtStringCiHi}{30.12}% data/runs/2026-09-17T08-42-17Z-keypath-mechanism/keypath_stats.json
\newcommand{\HostJwtStringCiLo}{29.75}% data/runs/2026-09-17T08-42-17Z-keypath-mechanism/keypath_stats.json
\newcommand{\HostJwtStringSafe}{7.208}% data/runs/2026-09-17T08-42-17Z-keypath-mechanism/keypath_stats.json
\newcommand{\HostJwtStringSafeCiHi}{7.334}% data/runs/2026-09-17T08-42-17Z-keypath-mechanism/keypath_stats.json
\newcommand{\HostJwtStringSafeCiLo}{7.125}% data/runs/2026-09-17T08-42-17Z-keypath-mechanism/keypath_stats.json
\newcommand{\HostNodeVersion}{26.6.0}% data/runs/2026-09-17T08-42-17Z-keypath-mechanism/keypath_stats.json
\newcommand{\HostOpensslAfterDrift}{3.6.4}% data/runs/HOST-OPENSSL-DRIFT-2026-09-17.md
\newcommand{\HostOpensslInferred}{3.6.3}% data/runs/HOST-OPENSSL-DRIFT-2026-09-17.md
\newcommand{\HostProbeShareOfCall}{60}% data/runs/2026-09-17T08-42-17Z-keypath-mechanism/keypath_stats.json
\newcommand{\HostProbeSucceeds}{13.83}% data/runs/2026-09-17T08-42-17Z-keypath-mechanism/keypath_stats.json
\newcommand{\HostProbeSucceedsCiHi}{13.92}% data/runs/2026-09-17T08-42-17Z-keypath-mechanism/keypath_stats.json
\newcommand{\HostProbeSucceedsCiLo}{13.62}% data/runs/2026-09-17T08-42-17Z-keypath-mechanism/keypath_stats.json
\newcommand{\HostProbeThrows}{18.04}% data/runs/2026-09-17T08-42-17Z-keypath-mechanism/keypath_stats.json
\newcommand{\HostProbeThrowsCiHi}{18.25}% data/runs/2026-09-17T08-42-17Z-keypath-mechanism/keypath_stats.json
\newcommand{\HostProbeThrowsCiLo}{17.83}% data/runs/2026-09-17T08-42-17Z-keypath-mechanism/keypath_stats.json
\newcommand{\HostRsStringPenalty}{27.29}% data/runs/2026-09-17T08-42-17Z-keypath-mechanism/keypath_stats.json
\newcommand{\HostRsStringPenaltyCiHi}{27.37}% data/runs/2026-09-17T08-42-17Z-keypath-mechanism/keypath_stats.json
\newcommand{\HostRsStringPenaltyCiLo}{27.00}% data/runs/2026-09-17T08-42-17Z-keypath-mechanism/keypath_stats.json
\newcommand{\HostRsStringPenaltyRounds}{15}% data/runs/2026-09-17T08-42-17Z-keypath-mechanism/keypath_stats.json
\newcommand{\HostSafePatchRatio}{4.15}% data/runs/2026-09-17T08-42-17Z-keypath-mechanism/keypath_stats.json
\newcommand{\HostSafePatchResidual}{1.208}% data/runs/2026-09-17T08-42-17Z-keypath-mechanism/keypath_stats.json
\newcommand{\HostSafePatchResidualCiHi}{1.251}% data/runs/2026-09-17T08-42-17Z-keypath-mechanism/keypath_stats.json
\newcommand{\HostSafePatchResidualCiLo}{1.084}% data/runs/2026-09-17T08-42-17Z-keypath-mechanism/keypath_stats.json
\newcommand{\HostSafePatchResidualRounds}{15}% data/runs/2026-09-17T08-42-17Z-keypath-mechanism/keypath_stats.json
\newcommand{\HostSafePatchSaving}{22.71}% data/runs/2026-09-17T08-42-17Z-keypath-mechanism/keypath_stats.json
\newcommand{\HostSafePatchSavingCiHi}{22.88}% data/runs/2026-09-17T08-42-17Z-keypath-mechanism/keypath_stats.json
\newcommand{\HostSafePatchSavingCiLo}{22.54}% data/runs/2026-09-17T08-42-17Z-keypath-mechanism/keypath_stats.json
\newcommand{\HostSafePatchSavingRounds}{15}% data/runs/2026-09-17T08-42-17Z-keypath-mechanism/keypath_stats.json
\newcommand{\HostStringPenalty}{23.92}% data/runs/2026-09-17T08-42-17Z-keypath-mechanism/keypath_stats.json
\newcommand{\HostStringPenaltyCiHi}{24.00}% data/runs/2026-09-17T08-42-17Z-keypath-mechanism/keypath_stats.json
\newcommand{\HostStringPenaltyCiLo}{23.58}% data/runs/2026-09-17T08-42-17Z-keypath-mechanism/keypath_stats.json
\newcommand{\HostStringPenaltyRounds}{15}% data/runs/2026-09-17T08-42-17Z-keypath-mechanism/keypath_stats.json
\newcommand{\HostTimerOverhead}{0.042}% data/runs/2026-09-17T08-42-17Z-keypath-mechanism/keypath_stats.json
\newcommand{\HostTimerOverheadCiHi}{0.042}% data/runs/2026-09-17T08-42-17Z-keypath-mechanism/keypath_stats.json
\newcommand{\HostTimerOverheadCiLo}{0.042}% data/runs/2026-09-17T08-42-17Z-keypath-mechanism/keypath_stats.json
\newcommand{\HostVeight}{14}% data/runs/2026-09-17T08-42-17Z-keypath-mechanism/run_metadata.json
\newcommand{\HostVeightFull}{14.6.202.34-node.26}% data/runs/2026-09-17T08-42-17Z-keypath-mechanism/run_metadata.json
\newcommand{\IdentityMismatches}{0}% data/runs/2026-09-17T08-57-13Z-keypath-runtime-matrix/environments.csv
\newcommand{\InsituObservedHi}{1028}% data/runs/2026-09-14T05-57-22Z-verifyrate-hs256/verify_rate_curve.csv
\newcommand{\InsituObservedLo}{453}% data/runs/2026-09-14T05-57-22Z-verifyrate-hs256/verify_rate_curve.csv
\newcommand{\InsituPredictedFromMatrix}{381}% data/runs/2026-09-17T08-57-13Z-keypath-runtime-matrix/keypath_stats.json
\newcommand{\InsituPreparsedAscA}{221.42}% data/runs/2026-09-14T15-01-36Z-verifyrate-hs256-preparsed/verify_rate_curve.csv
\newcommand{\InsituPreparsedAscB}{235.01}% data/runs/2026-09-14T15-01-36Z-verifyrate-hs256-preparsed/verify_rate_curve.csv
\newcommand{\InsituPreparsedAscC}{168.44}% data/runs/2026-09-14T15-01-36Z-verifyrate-hs256-preparsed/verify_rate_curve.csv
\newcommand{\InsituPreparsedAscD}{127.19}% data/runs/2026-09-14T15-01-36Z-verifyrate-hs256-preparsed/verify_rate_curve.csv
\newcommand{\InsituPreparsedAscE}{79.70}% data/runs/2026-09-14T15-01-36Z-verifyrate-hs256-preparsed/verify_rate_curve.csv
\newcommand{\InsituPreparsedAscF}{71.68}% data/runs/2026-09-14T15-01-36Z-verifyrate-hs256-preparsed/verify_rate_curve.csv
\newcommand{\InsituPreparsedAscG}{39.38}% data/runs/2026-09-14T15-01-36Z-verifyrate-hs256-preparsed/verify_rate_curve.csv
\newcommand{\InsituPreparsedDescA}{200.88}% data/runs/2026-09-14T15-01-36Z-verifyrate-hs256-preparsed/verify_rate_curve.csv
\newcommand{\InsituPreparsedDescB}{188.06}% data/runs/2026-09-14T15-01-36Z-verifyrate-hs256-preparsed/verify_rate_curve.csv
\newcommand{\InsituPreparsedDescC}{160.58}% data/runs/2026-09-14T15-01-36Z-verifyrate-hs256-preparsed/verify_rate_curve.csv
\newcommand{\InsituPreparsedDescD}{116.67}% data/runs/2026-09-14T15-01-36Z-verifyrate-hs256-preparsed/verify_rate_curve.csv
\newcommand{\InsituPreparsedDescE}{66.93}% data/runs/2026-09-14T15-01-36Z-verifyrate-hs256-preparsed/verify_rate_curve.csv
\newcommand{\InsituPreparsedDescF}{42.78}% data/runs/2026-09-14T15-01-36Z-verifyrate-hs256-preparsed/verify_rate_curve.csv
\newcommand{\InsituPreparsedDescG}{34.61}% data/runs/2026-09-14T15-01-36Z-verifyrate-hs256-preparsed/verify_rate_curve.csv
\newcommand{\InsituPreparsedPoints}{7}% data/runs/2026-09-14T15-01-36Z-verifyrate-hs256-preparsed/verify_rate_curve.csv
\newcommand{\InsituPreparsedSameSign}{7}% data/runs/2026-09-14T15-01-36Z-verifyrate-hs256-preparsed/verify_rate_curve.csv
\newcommand{\InsituRateA}{5}% data/runs/2026-09-14T05-57-22Z-verifyrate-hs256/verify_rate_curve.csv
\newcommand{\InsituRateB}{10}% data/runs/2026-09-14T05-57-22Z-verifyrate-hs256/verify_rate_curve.csv
\newcommand{\InsituRateC}{25}% data/runs/2026-09-14T05-57-22Z-verifyrate-hs256/verify_rate_curve.csv
\newcommand{\InsituRateD}{50}% data/runs/2026-09-14T05-57-22Z-verifyrate-hs256/verify_rate_curve.csv
\newcommand{\InsituRateE}{100}% data/runs/2026-09-14T05-57-22Z-verifyrate-hs256/verify_rate_curve.csv
\newcommand{\InsituRateF}{200}% data/runs/2026-09-14T05-57-22Z-verifyrate-hs256/verify_rate_curve.csv
\newcommand{\InsituRateG}{400}% data/runs/2026-09-14T05-57-22Z-verifyrate-hs256/verify_rate_curve.csv
\newcommand{\InsituRsAscA}{351.05}% data/runs/2026-09-14T06-41-56Z-verifyrate-rs256/verify_rate_curve.csv
\newcommand{\InsituRsAscB}{338.89}% data/runs/2026-09-14T06-41-56Z-verifyrate-rs256/verify_rate_curve.csv
\newcommand{\InsituRsAscC}{270.77}% data/runs/2026-09-14T06-41-56Z-verifyrate-rs256/verify_rate_curve.csv
\newcommand{\InsituRsAscD}{203.19}% data/runs/2026-09-14T06-41-56Z-verifyrate-rs256/verify_rate_curve.csv
\newcommand{\InsituRsAscE}{120.42}% data/runs/2026-09-14T06-41-56Z-verifyrate-rs256/verify_rate_curve.csv
\newcommand{\InsituRsAscF}{73.92}% data/runs/2026-09-14T06-41-56Z-verifyrate-rs256/verify_rate_curve.csv
\newcommand{\InsituRsAscG}{59.11}% data/runs/2026-09-14T06-41-56Z-verifyrate-rs256/verify_rate_curve.csv
\newcommand{\InsituRsDescA}{292.11}% data/runs/2026-09-14T06-41-56Z-verifyrate-rs256/verify_rate_curve.csv
\newcommand{\InsituRsDescB}{286.13}% data/runs/2026-09-14T06-41-56Z-verifyrate-rs256/verify_rate_curve.csv
\newcommand{\InsituRsDescC}{251.70}% data/runs/2026-09-14T06-41-56Z-verifyrate-rs256/verify_rate_curve.csv
\newcommand{\InsituRsDescD}{193.82}% data/runs/2026-09-14T06-41-56Z-verifyrate-rs256/verify_rate_curve.csv
\newcommand{\InsituRsDescE}{111.88}% data/runs/2026-09-14T06-41-56Z-verifyrate-rs256/verify_rate_curve.csv
\newcommand{\InsituRsDescF}{73.94}% data/runs/2026-09-14T06-41-56Z-verifyrate-rs256/verify_rate_curve.csv
\newcommand{\InsituRsDescG}{58.88}% data/runs/2026-09-14T06-41-56Z-verifyrate-rs256/verify_rate_curve.csv
\newcommand{\InsituRsPoints}{7}% data/runs/2026-09-14T06-41-56Z-verifyrate-rs256/verify_rate_curve.csv
\newcommand{\InsituRsSameSign}{6}% data/runs/2026-09-14T06-41-56Z-verifyrate-rs256/verify_rate_curve.csv
\newcommand{\InsituStringAscA}{928.24}% data/runs/2026-09-14T05-57-22Z-verifyrate-hs256/verify_rate_curve.csv
\newcommand{\InsituStringAscB}{976.28}% data/runs/2026-09-14T05-57-22Z-verifyrate-hs256/verify_rate_curve.csv
\newcommand{\InsituStringAscC}{918.54}% data/runs/2026-09-14T05-57-22Z-verifyrate-hs256/verify_rate_curve.csv
\newcommand{\InsituStringAscD}{758.57}% data/runs/2026-09-14T05-57-22Z-verifyrate-hs256/verify_rate_curve.csv
\newcommand{\InsituStringAscE}{571.59}% data/runs/2026-09-14T05-57-22Z-verifyrate-hs256/verify_rate_curve.csv
\newcommand{\InsituStringAscF}{494.91}% data/runs/2026-09-14T05-57-22Z-verifyrate-hs256/verify_rate_curve.csv
\newcommand{\InsituStringAscG}{452.82}% data/runs/2026-09-14T05-57-22Z-verifyrate-hs256/verify_rate_curve.csv
\newcommand{\InsituStringDescA}{1027.97}% data/runs/2026-09-14T05-57-22Z-verifyrate-hs256/verify_rate_curve.csv
\newcommand{\InsituStringDescB}{1017.34}% data/runs/2026-09-14T05-57-22Z-verifyrate-hs256/verify_rate_curve.csv
\newcommand{\InsituStringDescC}{934.47}% data/runs/2026-09-14T05-57-22Z-verifyrate-hs256/verify_rate_curve.csv
\newcommand{\InsituStringDescD}{737.12}% data/runs/2026-09-14T05-57-22Z-verifyrate-hs256/verify_rate_curve.csv
\newcommand{\InsituStringDescE}{607.78}% data/runs/2026-09-14T05-57-22Z-verifyrate-hs256/verify_rate_curve.csv
\newcommand{\InsituStringDescF}{532.50}% data/runs/2026-09-14T05-57-22Z-verifyrate-hs256/verify_rate_curve.csv
\newcommand{\InsituStringDescG}{467.62}% data/runs/2026-09-14T05-57-22Z-verifyrate-hs256/verify_rate_curve.csv
\newcommand{\InsituStringPoints}{7}% data/runs/2026-09-14T05-57-22Z-verifyrate-hs256/verify_rate_curve.csv
\newcommand{\InsituStringSameSign}{6}% data/runs/2026-09-14T05-57-22Z-verifyrate-hs256/verify_rate_curve.csv
\newcommand{\InsituTightestAgreement}{0.02}% data/runs/2026-09-14T06-41-56Z-verifyrate-rs256/verify_rate_curve.csv
\newcommand{\InsituUnaccountedHi}{647}% data/runs/2026-09-14T05-57-22Z-verifyrate-hs256/verify_rate_curve.csv
\newcommand{\InsituUnaccountedLo}{71}% data/runs/2026-09-14T05-57-22Z-verifyrate-hs256/verify_rate_curve.csv
\newcommand{\ManifestFiles}{502}% survey/data/corpus_manifest.csv
\newcommand{\ManifestHeadResolved}{501}% survey/data/corpus_manifest.csv
\newcommand{\ManifestRepos}{457}% survey/data/corpus_manifest.csv
\newcommand{\ManifestSampleFiles}{311}% survey/data/corpus_manifest.csv
\newcommand{\MatrixCallsPerInvocation}{20000}% data/runs/2026-09-17T08-57-13Z-keypath-runtime-matrix/keypath_stats.json
\newcommand{\MatrixCvMax}{20.2}% data/runs/2026-09-17T08-57-13Z-keypath-runtime-matrix/keypath_stats.json
\newcommand{\MatrixCvMaxCondition}{probe\_succeeds}% data/runs/2026-09-17T08-57-13Z-keypath-runtime-matrix/keypath_stats.json
\newcommand{\MatrixCvMaxEnvironment}{node:26.6.0-bookworm}% data/runs/2026-09-17T08-57-13Z-keypath-runtime-matrix/keypath_stats.json
\newcommand{\MatrixCvRatio}{5.7}% data/runs/2026-09-17T08-57-13Z-keypath-runtime-matrix/keypath_stats.json
\newcommand{\MatrixInvocations}{10}% data/runs/2026-09-17T08-57-13Z-keypath-runtime-matrix/keypath_stats.json
\newcommand{\NodeEighteenDecodeOnly}{1.62}% data/runs/2026-09-17T08-57-13Z-keypath-runtime-matrix/keypath_stats.json
\newcommand{\NodeEighteenHmacString}{2.42}% data/runs/2026-09-17T08-57-13Z-keypath-runtime-matrix/keypath_stats.json
\newcommand{\NodeEighteenImageId}{8d6421d663b4}% data/runs/runtime_identity.csv
\newcommand{\NodeEighteenJwtPreparsed}{7.58}% data/runs/2026-09-17T08-57-13Z-keypath-runtime-matrix/keypath_stats.json
\newcommand{\NodeEighteenJwtString}{394.71}% data/runs/2026-09-17T08-57-13Z-keypath-runtime-matrix/keypath_stats.json
\newcommand{\NodeEighteenNodeLiteral}{18.20.8}% data/runs/2026-09-17T08-57-13Z-keypath-runtime-matrix/keypath_stats.json
\newcommand{\NodeEighteenNodeMajor}{18}% data/runs/2026-09-17T08-57-13Z-keypath-runtime-matrix/keypath_stats.json
\newcommand{\NodeEighteenOpensslLiteral}{3.0.16}% data/runs/2026-09-17T08-57-13Z-keypath-runtime-matrix/keypath_stats.json
\newcommand{\NodeEighteenOpensslSeries}{3.0}% data/runs/2026-09-17T08-57-13Z-keypath-runtime-matrix/keypath_stats.json
\newcommand{\NodeEighteenProbeSucceeds}{115.83}% data/runs/2026-09-17T08-57-13Z-keypath-runtime-matrix/keypath_stats.json
\newcommand{\NodeEighteenProbeThrows}{373.79}% data/runs/2026-09-17T08-57-13Z-keypath-runtime-matrix/keypath_stats.json
\newcommand{\NodeEighteenProbeThrowsCiHi}{375.79}% data/runs/2026-09-17T08-57-13Z-keypath-runtime-matrix/keypath_stats.json
\newcommand{\NodeEighteenProbeThrowsCiLo}{371.50}% data/runs/2026-09-17T08-57-13Z-keypath-runtime-matrix/keypath_stats.json
\newcommand{\NodeEighteenVeight}{10}% data/runs/runtime_identity.csv
\newcommand{\NodeEighteenVeightFull}{10.2.154.26-node.39}% data/runs/runtime_identity.csv
\newcommand{\NodeTwentyDecodeOnly}{2.19}% data/runs/2026-09-17T08-57-13Z-keypath-runtime-matrix/keypath_stats.json
\newcommand{\NodeTwentyHmacString}{2.33}% data/runs/2026-09-17T08-57-13Z-keypath-runtime-matrix/keypath_stats.json
\newcommand{\NodeTwentyImageId}{fb4cd12c85ee}% data/runs/runtime_identity.csv
\newcommand{\NodeTwentyJwtPreparsed}{7.25}% data/runs/2026-09-17T08-57-13Z-keypath-runtime-matrix/keypath_stats.json
\newcommand{\NodeTwentyJwtString}{437.40}% data/runs/2026-09-17T08-57-13Z-keypath-runtime-matrix/keypath_stats.json
\newcommand{\NodeTwentyNodeLiteral}{20.20.2}% data/runs/2026-09-17T08-57-13Z-keypath-runtime-matrix/keypath_stats.json
\newcommand{\NodeTwentyNodeMajor}{20}% data/runs/2026-09-17T08-57-13Z-keypath-runtime-matrix/keypath_stats.json
\newcommand{\NodeTwentyOpensslLiteral}{3.0.19}% data/runs/2026-09-17T08-57-13Z-keypath-runtime-matrix/keypath_stats.json
\newcommand{\NodeTwentyOpensslSeries}{3.0}% data/runs/2026-09-17T08-57-13Z-keypath-runtime-matrix/keypath_stats.json
\newcommand{\NodeTwentyProbeSucceeds}{142.38}% data/runs/2026-09-17T08-57-13Z-keypath-runtime-matrix/keypath_stats.json
\newcommand{\NodeTwentyProbeThrows}{420.62}% data/runs/2026-09-17T08-57-13Z-keypath-runtime-matrix/keypath_stats.json
\newcommand{\NodeTwentyProbeThrowsCiHi}{422.56}% data/runs/2026-09-17T08-57-13Z-keypath-runtime-matrix/keypath_stats.json
\newcommand{\NodeTwentyProbeThrowsCiLo}{419.52}% data/runs/2026-09-17T08-57-13Z-keypath-runtime-matrix/keypath_stats.json
\newcommand{\NodeTwentySixGlibcDecodeOnly}{1.67}% data/runs/2026-09-17T08-57-13Z-keypath-runtime-matrix/keypath_stats.json
\newcommand{\NodeTwentySixGlibcHmacString}{3.19}% data/runs/2026-09-17T08-57-13Z-keypath-runtime-matrix/keypath_stats.json
\newcommand{\NodeTwentySixGlibcImageId}{6e0a69f43d2b}% data/runs/runtime_identity.csv
\newcommand{\NodeTwentySixGlibcJwtPreparsed}{6.04}% data/runs/2026-09-17T08-57-13Z-keypath-runtime-matrix/keypath_stats.json
\newcommand{\NodeTwentySixGlibcJwtString}{27.96}% data/runs/2026-09-17T08-57-13Z-keypath-runtime-matrix/keypath_stats.json
\newcommand{\NodeTwentySixGlibcNodeLiteral}{26.6.0}% data/runs/2026-09-17T08-57-13Z-keypath-runtime-matrix/keypath_stats.json
\newcommand{\NodeTwentySixGlibcNodeMajor}{26}% data/runs/2026-09-17T08-57-13Z-keypath-runtime-matrix/keypath_stats.json
\newcommand{\NodeTwentySixGlibcOpensslLiteral}{3.5.7}% data/runs/2026-09-17T08-57-13Z-keypath-runtime-matrix/keypath_stats.json
\newcommand{\NodeTwentySixGlibcOpensslSeries}{3.5}% data/runs/2026-09-17T08-57-13Z-keypath-runtime-matrix/keypath_stats.json
\newcommand{\NodeTwentySixGlibcProbeSucceeds}{12.12}% data/runs/2026-09-17T08-57-13Z-keypath-runtime-matrix/keypath_stats.json
\newcommand{\NodeTwentySixGlibcProbeThrows}{16.58}% data/runs/2026-09-17T08-57-13Z-keypath-runtime-matrix/keypath_stats.json
\newcommand{\NodeTwentySixGlibcProbeThrowsCiHi}{17.88}% data/runs/2026-09-17T08-57-13Z-keypath-runtime-matrix/keypath_stats.json
\newcommand{\NodeTwentySixGlibcProbeThrowsCiLo}{16.38}% data/runs/2026-09-17T08-57-13Z-keypath-runtime-matrix/keypath_stats.json
\newcommand{\NodeTwentySixGlibcVeight}{14}% data/runs/runtime_identity.csv
\newcommand{\NodeTwentySixGlibcVeightFull}{14.6.202.34-node.26}% data/runs/runtime_identity.csv
\newcommand{\NodeTwentySixMuslDecodeOnly}{1.71}% data/runs/2026-09-17T08-57-13Z-keypath-runtime-matrix/keypath_stats.json
\newcommand{\NodeTwentySixMuslHmacString}{3.25}% data/runs/2026-09-17T08-57-13Z-keypath-runtime-matrix/keypath_stats.json
\newcommand{\NodeTwentySixMuslImageId}{a4fb14143ee2}% data/runs/runtime_identity.csv
\newcommand{\NodeTwentySixMuslJwtPreparsed}{6.56}% data/runs/2026-09-17T08-57-13Z-keypath-runtime-matrix/keypath_stats.json
\newcommand{\NodeTwentySixMuslJwtString}{36.42}% data/runs/2026-09-17T08-57-13Z-keypath-runtime-matrix/keypath_stats.json
\newcommand{\NodeTwentySixMuslNodeLiteral}{26.6.0}% data/runs/2026-09-17T08-57-13Z-keypath-runtime-matrix/keypath_stats.json
\newcommand{\NodeTwentySixMuslNodeMajor}{26}% data/runs/2026-09-17T08-57-13Z-keypath-runtime-matrix/keypath_stats.json
\newcommand{\NodeTwentySixMuslOpensslLiteral}{3.5.7}% data/runs/2026-09-17T08-57-13Z-keypath-runtime-matrix/keypath_stats.json
\newcommand{\NodeTwentySixMuslOpensslSeries}{3.5}% data/runs/2026-09-17T08-57-13Z-keypath-runtime-matrix/keypath_stats.json
\newcommand{\NodeTwentySixMuslProbeSucceeds}{15.54}% data/runs/2026-09-17T08-57-13Z-keypath-runtime-matrix/keypath_stats.json
\newcommand{\NodeTwentySixMuslProbeThrows}{22.94}% data/runs/2026-09-17T08-57-13Z-keypath-runtime-matrix/keypath_stats.json
\newcommand{\NodeTwentySixMuslProbeThrowsCiHi}{23.29}% data/runs/2026-09-17T08-57-13Z-keypath-runtime-matrix/keypath_stats.json
\newcommand{\NodeTwentySixMuslProbeThrowsCiLo}{22.79}% data/runs/2026-09-17T08-57-13Z-keypath-runtime-matrix/keypath_stats.json
\newcommand{\NodeTwentySixMuslVeight}{14}% data/runs/runtime_identity.csv
\newcommand{\NodeTwentySixMuslVeightFull}{14.6.202.34-node.26}% data/runs/runtime_identity.csv
\newcommand{\NodeTwentyVeight}{11}% data/runs/runtime_identity.csv
\newcommand{\NodeTwentyVeightFull}{11.3.244.8-node.38}% data/runs/runtime_identity.csv
\newcommand{\PopBestCellPct}{0.0000}% survey/data/population_counts.csv
\newcommand{\PopCapturedAt}{2026-09-18}% survey/data/population_counts.csv
\newcommand{\PopJsImportFiles}{524288}% survey/data/population_counts.csv
\newcommand{\PopJsImportFilesWithApi}{38}% survey/data/population_counts.csv
\newcommand{\PopJsImportPct}{0.0072}% survey/data/population_counts.csv
\newcommand{\PopJsRequireFiles}{790528}% survey/data/population_counts.csv
\newcommand{\PopJsRequireFilesWithApi}{29}% survey/data/population_counts.csv
\newcommand{\PopJsRequirePct}{0.0037}% survey/data/population_counts.csv
\newcommand{\PopLargestCellFiles}{790528}% survey/data/population_counts.csv
\newcommand{\PopSumFiles}{1656928}% survey/data/population_counts.csv
\newcommand{\PopSumFilesWithApi}{168}% survey/data/population_counts.csv
\newcommand{\PopSumOneIn}{9863}% survey/data/population_counts.csv
\newcommand{\PopSumPct}{0.0101}% survey/data/population_counts.csv
\newcommand{\PopTsImportFiles}{333824}% survey/data/population_counts.csv
\newcommand{\PopTsImportFilesWithApi}{101}% survey/data/population_counts.csv
\newcommand{\PopTsImportPct}{0.0303}% survey/data/population_counts.csv
\newcommand{\PopTsRequireFiles}{8288}% survey/data/population_counts.csv
\newcommand{\PopTsRequireFilesWithApi}{0}% survey/data/population_counts.csv
\newcommand{\PopTsRequirePct}{0.0000}% survey/data/population_counts.csv
\newcommand{\PopWorstCellPct}{0.0303}% survey/data/population_counts.csv
\newcommand{\PreparsedUnattributed}{2.58}% data/runs/2026-09-17T08-42-17Z-keypath-mechanism/keypath_stats.json
\newcommand{\ProfEighteenPct}{95.3}% data/runs/2026-09-18T-profile-matched/cpuprof_matched.csv
\newcommand{\ProfEighteenPctHi}{95.6}% data/runs/2026-09-18T-profile-matched/cpuprof_matched.csv
\newcommand{\ProfEighteenPctLo}{95.1}% data/runs/2026-09-18T-profile-matched/cpuprof_matched.csv
\newcommand{\ProfHostPct}{66.2}% data/runs/2026-09-18T-profile-matched/cpuprof_matched.csv
\newcommand{\ProfHostPctHi}{67.4}% data/runs/2026-09-18T-profile-matched/cpuprof_matched.csv
\newcommand{\ProfHostPctLo}{65.3}% data/runs/2026-09-18T-profile-matched/cpuprof_matched.csv
\newcommand{\ProfIterations}{50000}% data/runs/2026-09-18T-profile-matched/cpuprof_matched.csv
\newcommand{\ProfRepeats}{3}% data/runs/2026-09-18T-profile-matched/cpuprof_matched.csv
\newcommand{\ProfTwentySixPct}{61.5}% data/runs/2026-09-18T-profile-matched/cpuprof_matched.csv
\newcommand{\ProfTwentySixPctHi}{61.6}% data/runs/2026-09-18T-profile-matched/cpuprof_matched.csv
\newcommand{\ProfTwentySixPctLo}{58.5}% data/runs/2026-09-18T-profile-matched/cpuprof_matched.csv
\newcommand{\RematchAdjudicated}{12}% survey/data/hand_adjudication.csv
\newcommand{\RematchSitesPreparsed}{0}% survey/data/hand_adjudication.csv
\newcommand{\SampleBuffer}{2}% survey/data/callsites.csv
\newcommand{\SampleCallSites}{327}% survey/data/callsites.csv
\newcommand{\SampleEnvString}{176}% survey/data/callsites.csv
\newcommand{\SampleFileContents}{4}% survey/data/callsites.csv
\newcommand{\SampleNoAlgorithms}{309}% survey/data/callsites.csv
\newcommand{\SampleNoAlgorithmsPct}{94}% survey/data/callsites.csv
\newcommand{\SampleRematchAdded}{12}% survey/data/callsites.csv
\newcommand{\SampleRepos}{247}% survey/data/callsites.csv
\newcommand{\SampleSitesPreparsed}{0}% survey/data/callsites.csv
\newcommand{\SampleStringLiteral}{52}% survey/data/callsites.csv
\newcommand{\SampleUnknown}{93}% survey/data/callsites.csv
\newcommand{\SpanDecodeOnly}{1.35}% data/runs/2026-09-17T08-57-13Z-keypath-runtime-matrix/keypath_stats.json
\newcommand{\SpanDecodeOnlyWithHost}{1.69}% data/runs/2026-09-17T08-57-13Z-keypath-runtime-matrix/keypath_stats.json
\newcommand{\SpanHmacString}{1.39}% data/runs/2026-09-17T08-57-13Z-keypath-runtime-matrix/keypath_stats.json
\newcommand{\SpanHmacStringWithHost}{1.47}% data/runs/2026-09-17T08-57-13Z-keypath-runtime-matrix/keypath_stats.json
\newcommand{\SpanJwtPreparsed}{1.26}% data/runs/2026-09-17T08-57-13Z-keypath-runtime-matrix/keypath_stats.json
\newcommand{\SpanJwtPreparsedWithHost}{1.26}% data/runs/2026-09-17T08-57-13Z-keypath-runtime-matrix/keypath_stats.json
\newcommand{\SpanProbeSucceeds}{11.74}% data/runs/2026-09-17T08-57-13Z-keypath-runtime-matrix/keypath_stats.json
\newcommand{\SpanProbeSucceedsWithHost}{11.74}% data/runs/2026-09-17T08-57-13Z-keypath-runtime-matrix/keypath_stats.json
\newcommand{\SpanProbeThrows}{25.36}% data/runs/2026-09-17T08-57-13Z-keypath-runtime-matrix/keypath_stats.json
\newcommand{\SpanProbeThrowsWithHost}{25.36}% data/runs/2026-09-17T08-57-13Z-keypath-runtime-matrix/keypath_stats.json
\newcommand{\SuiteNaiveFailing}{2}% data/runs/2026-09-18T-upstream-suite/suite_results.csv
\newcommand{\SuiteNaivePassing}{509}% data/runs/2026-09-18T-upstream-suite/suite_results.csv
\newcommand{\SuitePatchedFailing}{0}% data/runs/2026-09-18T-upstream-suite/suite_results.csv
\newcommand{\SuitePatchedPassing}{511}% data/runs/2026-09-18T-upstream-suite/suite_results.csv
\newcommand{\SuiteStockFailing}{0}% data/runs/2026-09-18T-upstream-suite/suite_results.csv
\newcommand{\SuiteStockPassing}{511}% data/runs/2026-09-18T-upstream-suite/suite_results.csv

% --- declared but unmeasured ------------------------------------------
% xEightSixReplication: x86_64 replication of the runtime matrix. Not performed: the Oracle Cloud instance was unreachable. See threats to validity.
\newcommand{\xEightSixReplication}{\PLACEHOLDER{xEightSixReplication}}

% Generated by analysis/make_revision_macros.py -- do not edit by hand.
% Each macro is read from the file named in revision_macros_provenance.csv.
\providecommand{\PLACEHOLDER}[1]{\textcolor{red}{\textbf{[#1]}}}
\newcommand{\CErrorStringsVThreeFiveEight}{0.483}
\newcommand{\CErrorStringsVThreeFiveEightCiHi}{0.491}
\newcommand{\CErrorStringsVThreeFiveEightCiLo}{0.479}
\newcommand{\CErrorStringsVThreeFourThree}{0.500}
\newcommand{\CErrorStringsVThreeFourThreeCiHi}{0.503}
\newcommand{\CErrorStringsVThreeFourThreeCiLo}{0.495}
\newcommand{\CErrorStringsVThreeOneEight}{0.497}
\newcommand{\CErrorStringsVThreeOneEightCiHi}{0.500}
\newcommand{\CErrorStringsVThreeOneEightCiLo}{0.493}
\newcommand{\CErrorStringsVThreeSixOne}{0.472}
\newcommand{\CErrorStringsVThreeSixOneCiHi}{0.475}
\newcommand{\CErrorStringsVThreeSixOneCiLo}{0.470}
\newcommand{\CErrorStringsVThreeThreeFive}{0.474}
\newcommand{\CErrorStringsVThreeThreeFiveCiHi}{0.475}
\newcommand{\CErrorStringsVThreeThreeFiveCiLo}{0.473}
\newcommand{\CErrorStringsVThreeTwoSix}{0.459}
\newcommand{\CErrorStringsVThreeTwoSixCiHi}{0.460}
\newcommand{\CErrorStringsVThreeTwoSixCiLo}{0.457}
\newcommand{\CErrorStringsVThreeZeroNineteen}{0.869}
\newcommand{\CErrorStringsVThreeZeroNineteenCiHi}{0.874}
\newcommand{\CErrorStringsVThreeZeroNineteenCiLo}{0.859}
\newcommand{\CErrorStringsVThreeZeroSixteen}{0.796}
\newcommand{\CErrorStringsVThreeZeroSixteenCiHi}{0.818}
\newcommand{\CErrorStringsVThreeZeroSixteenCiLo}{0.659}
\newcommand{\CFirstFastVersion}{3.1.8}
\newcommand{\CFullCiMaxHalfWidthPct}{0.9}
\newcommand{\CFullRatioOldNew}{46.69}
\newcommand{\CFullRatioOldNewCiHi}{46.90}
\newcommand{\CFullRatioOldNewCiLo}{46.16}
\newcommand{\CFullSpan}{62.10}
\newcommand{\CFullVThreeFiveEight}{8.447}
\newcommand{\CFullVThreeFiveEightCiHi}{8.475}
\newcommand{\CFullVThreeFiveEightCiLo}{8.432}
\newcommand{\CFullVThreeFourThree}{7.245}
\newcommand{\CFullVThreeFourThreeCiHi}{7.276}
\newcommand{\CFullVThreeFourThreeCiLo}{7.220}
\newcommand{\CFullVThreeOneEight}{57.48}
\newcommand{\CFullVThreeOneEightCiHi}{57.73}
\newcommand{\CFullVThreeOneEightCiLo}{57.47}
\newcommand{\CFullVThreeSixOne}{8.803}
\newcommand{\CFullVThreeSixOneCiHi}{8.840}
\newcommand{\CFullVThreeSixOneCiLo}{8.726}
\newcommand{\CFullVThreeThreeFive}{6.946}
\newcommand{\CFullVThreeThreeFiveCiHi}{6.977}
\newcommand{\CFullVThreeThreeFiveCiLo}{6.914}
\newcommand{\CFullVThreeTwoSix}{6.998}
\newcommand{\CFullVThreeTwoSixCiHi}{7.030}
\newcommand{\CFullVThreeTwoSixCiLo}{6.983}
\newcommand{\CFullVThreeZeroNineteen}{431.3}
\newcommand{\CFullVThreeZeroNineteenCiHi}{431.8}
\newcommand{\CFullVThreeZeroNineteenCiLo}{427.9}
\newcommand{\CFullVThreeZeroSixteen}{394.4}
\newcommand{\CFullVThreeZeroSixteenCiHi}{396.5}
\newcommand{\CFullVThreeZeroSixteenCiLo}{390.9}
\newcommand{\CPrivateParseVThreeFiveEight}{5.757}
\newcommand{\CPrivateParseVThreeFiveEightCiHi}{5.771}
\newcommand{\CPrivateParseVThreeFiveEightCiLo}{5.742}
\newcommand{\CPrivateParseVThreeFourThree}{4.554}
\newcommand{\CPrivateParseVThreeFourThreeCiHi}{4.575}
\newcommand{\CPrivateParseVThreeFourThreeCiLo}{4.538}
\newcommand{\CPrivateParseVThreeOneEight}{54.84}
\newcommand{\CPrivateParseVThreeOneEightCiHi}{55.32}
\newcommand{\CPrivateParseVThreeOneEightCiLo}{54.81}
\newcommand{\CPrivateParseVThreeSixOne}{5.985}
\newcommand{\CPrivateParseVThreeSixOneCiHi}{6.026}
\newcommand{\CPrivateParseVThreeSixOneCiLo}{5.913}
\newcommand{\CPrivateParseVThreeThreeFive}{4.465}
\newcommand{\CPrivateParseVThreeThreeFiveCiHi}{4.496}
\newcommand{\CPrivateParseVThreeThreeFiveCiLo}{4.445}
\newcommand{\CPrivateParseVThreeTwoSix}{4.487}
\newcommand{\CPrivateParseVThreeTwoSixCiHi}{4.499}
\newcommand{\CPrivateParseVThreeTwoSixCiLo}{4.459}
\newcommand{\CPrivateParseVThreeZeroNineteen}{424.0}
\newcommand{\CPrivateParseVThreeZeroNineteenCiHi}{427.6}
\newcommand{\CPrivateParseVThreeZeroNineteenCiLo}{422.2}
\newcommand{\CPrivateParseVThreeZeroSixteen}{390.9}
\newcommand{\CPrivateParseVThreeZeroSixteenCiHi}{392.0}
\newcommand{\CPrivateParseVThreeZeroSixteenCiLo}{389.6}
\newcommand{\CPrivateShareVThreeFiveEight}{68}
\newcommand{\CPrivateShareVThreeFourThree}{63}
\newcommand{\CPrivateShareVThreeOneEight}{95}
\newcommand{\CPrivateShareVThreeSixOne}{68}
\newcommand{\CPrivateShareVThreeThreeFive}{64}
\newcommand{\CPrivateShareVThreeTwoSix}{64}
\newcommand{\CPrivateShareVThreeZeroNineteen}{98}
\newcommand{\CPrivateShareVThreeZeroSixteen}{99}
\newcommand{\CProbeCallsPerInvocation}{5000}
\newcommand{\CProbeCompiler}{GCC 13.3.0}
\newcommand{\CProbeCpuModel}{Intel(R) Xeon(R) Processor @ 2.80GHz}
\newcommand{\CProbeCvMax}{28.0}
\newcommand{\CProbeCvMedian}{1.5}
\newcommand{\CProbeErrorsQueued}{1}
\newcommand{\CProbeIdentityMismatches}{0}
\newcommand{\CProbeInvocations}{15}
\newcommand{\CProbeVersionCount}{8}
\newcommand{\CProbeVersionFirst}{3.0.16}
\newcommand{\CProbeVersionLast}{3.6.1}
\newcommand{\CPublicTriesVThreeFiveEight}{2.051}
\newcommand{\CPublicTriesVThreeFiveEightCiHi}{2.061}
\newcommand{\CPublicTriesVThreeFiveEightCiLo}{2.049}
\newcommand{\CPublicTriesVThreeFourThree}{2.098}
\newcommand{\CPublicTriesVThreeFourThreeCiHi}{2.101}
\newcommand{\CPublicTriesVThreeFourThreeCiLo}{2.092}
\newcommand{\CPublicTriesVThreeOneEight}{2.196}
\newcommand{\CPublicTriesVThreeOneEightCiHi}{2.205}
\newcommand{\CPublicTriesVThreeOneEightCiLo}{2.194}
\newcommand{\CPublicTriesVThreeSixOne}{2.234}
\newcommand{\CPublicTriesVThreeSixOneCiHi}{2.241}
\newcommand{\CPublicTriesVThreeSixOneCiLo}{2.230}
\newcommand{\CPublicTriesVThreeThreeFive}{2.074}
\newcommand{\CPublicTriesVThreeThreeFiveCiHi}{2.077}
\newcommand{\CPublicTriesVThreeThreeFiveCiLo}{2.071}
\newcommand{\CPublicTriesVThreeTwoSix}{2.057}
\newcommand{\CPublicTriesVThreeTwoSixCiHi}{2.067}
\newcommand{\CPublicTriesVThreeTwoSixCiLo}{2.056}
\newcommand{\CPublicTriesVThreeZeroNineteen}{2.231}
\newcommand{\CPublicTriesVThreeZeroNineteenCiHi}{2.239}
\newcommand{\CPublicTriesVThreeZeroNineteenCiLo}{2.228}
\newcommand{\CPublicTriesVThreeZeroSixteen}{2.287}
\newcommand{\CPublicTriesVThreeZeroSixteenCiHi}{2.294}
\newcommand{\CPublicTriesVThreeZeroSixteenCiLo}{2.284}
\newcommand{\CSpkiOkVThreeFiveEight}{11.26}
\newcommand{\CSpkiOkVThreeFiveEightCiHi}{11.29}
\newcommand{\CSpkiOkVThreeFiveEightCiLo}{11.23}
\newcommand{\CSpkiOkVThreeFourThree}{11.19}
\newcommand{\CSpkiOkVThreeFourThreeCiHi}{11.21}
\newcommand{\CSpkiOkVThreeFourThreeCiLo}{11.17}
\newcommand{\CSpkiOkVThreeOneEight}{36.09}
\newcommand{\CSpkiOkVThreeOneEightCiHi}{36.22}
\newcommand{\CSpkiOkVThreeOneEightCiLo}{36.03}
\newcommand{\CSpkiOkVThreeSixOne}{11.16}
\newcommand{\CSpkiOkVThreeSixOneCiHi}{11.18}
\newcommand{\CSpkiOkVThreeSixOneCiLo}{11.13}
\newcommand{\CSpkiOkVThreeThreeFive}{12.04}
\newcommand{\CSpkiOkVThreeThreeFiveCiHi}{12.09}
\newcommand{\CSpkiOkVThreeThreeFiveCiLo}{12.02}
\newcommand{\CSpkiOkVThreeTwoSix}{11.75}
\newcommand{\CSpkiOkVThreeTwoSixCiHi}{11.80}
\newcommand{\CSpkiOkVThreeTwoSixCiLo}{11.71}
\newcommand{\CSpkiOkVThreeZeroNineteen}{132.6}
\newcommand{\CSpkiOkVThreeZeroNineteenCiHi}{134.3}
\newcommand{\CSpkiOkVThreeZeroNineteenCiLo}{131.8}
\newcommand{\CSpkiOkVThreeZeroSixteen}{112.9}
\newcommand{\CSpkiOkVThreeZeroSixteenCiHi}{114.1}
\newcommand{\CSpkiOkVThreeZeroSixteenCiLo}{112.6}
\newcommand{\CTimerVThreeFiveEight}{0.026}
\newcommand{\CTimerVThreeFiveEightCiHi}{0.026}
\newcommand{\CTimerVThreeFiveEightCiLo}{0.026}
\newcommand{\CTimerVThreeFourThree}{0.026}
\newcommand{\CTimerVThreeFourThreeCiHi}{0.026}
\newcommand{\CTimerVThreeFourThreeCiLo}{0.026}
\newcommand{\CTimerVThreeOneEight}{0.026}
\newcommand{\CTimerVThreeOneEightCiHi}{0.026}
\newcommand{\CTimerVThreeOneEightCiLo}{0.026}
\newcommand{\CTimerVThreeSixOne}{0.026}
\newcommand{\CTimerVThreeSixOneCiHi}{0.026}
\newcommand{\CTimerVThreeSixOneCiLo}{0.026}
\newcommand{\CTimerVThreeThreeFive}{0.026}
\newcommand{\CTimerVThreeThreeFiveCiHi}{0.028}
\newcommand{\CTimerVThreeThreeFiveCiLo}{0.026}
\newcommand{\CTimerVThreeTwoSix}{0.026}
\newcommand{\CTimerVThreeTwoSixCiHi}{0.026}
\newcommand{\CTimerVThreeTwoSixCiLo}{0.026}
\newcommand{\CTimerVThreeZeroNineteen}{0.026}
\newcommand{\CTimerVThreeZeroNineteenCiHi}{0.026}
\newcommand{\CTimerVThreeZeroNineteenCiLo}{0.026}
\newcommand{\CTimerVThreeZeroSixteen}{0.026}
\newcommand{\CTimerVThreeZeroSixteenCiHi}{0.026}
\newcommand{\CTimerVThreeZeroSixteenCiLo}{0.026}
\newcommand{\CVersionLiteralVThreeFiveEight}{3.5.8}
\newcommand{\CVersionLiteralVThreeFourThree}{3.4.3}
\newcommand{\CVersionLiteralVThreeOneEight}{3.1.8}
\newcommand{\CVersionLiteralVThreeSixOne}{3.6.1}
\newcommand{\CVersionLiteralVThreeThreeFive}{3.3.5}
\newcommand{\CVersionLiteralVThreeTwoSix}{3.2.6}
\newcommand{\CVersionLiteralVThreeZeroNineteen}{3.0.19}
\newcommand{\CVersionLiteralVThreeZeroSixteen}{3.0.16}
\newcommand{\ExcCallsPerInvocation}{10000}
\newcommand{\ExcCpuModel}{Intel(R) Xeon(R) Processor @ 2.80GHz}
\newcommand{\ExcCvMax}{27.7}
\newcommand{\ExcCvMedian}{2.9}
\newcommand{\ExcDepth}{8}
\newcommand{\ExcErrorCode}{ERR\_OSSL\_UNSUPPORTED}
\newcommand{\ExcInvocations}{15}
\newcommand{\RtEighteenDecompC}{394.4}
\newcommand{\RtEighteenDecompCCiHi}{396.5}
\newcommand{\RtEighteenDecompCCiLo}{389.2}
\newcommand{\RtEighteenDecompCPct}{101}
\newcommand{\RtEighteenDecompCVersion}{3.0.16}
\newcommand{\RtEighteenDecompExcPct}{1.2}
\newcommand{\RtEighteenDecompLeftover}{-8.694}
\newcommand{\RtEighteenDecompLeftoverCiHi}{-3.497}
\newcommand{\RtEighteenDecompLeftoverCiLo}{-12.476}
\newcommand{\RtEighteenDecompLeftoverPct}{-2.2}
\newcommand{\RtEighteenExcNoStack}{2.204}
\newcommand{\RtEighteenExcNoStackCiHi}{2.232}
\newcommand{\RtEighteenExcNoStackCiLo}{2.200}
\newcommand{\RtEighteenExcNodeInternal}{12.88}
\newcommand{\RtEighteenExcNodeInternalCiHi}{13.23}
\newcommand{\RtEighteenExcNodeInternalCiLo}{12.71}
\newcommand{\RtEighteenExcNodeLike}{4.781}
\newcommand{\RtEighteenExcNodeLikeCiHi}{4.795}
\newcommand{\RtEighteenExcNodeLikeCiLo}{4.767}
\newcommand{\RtEighteenExcPlain}{4.867}
\newcommand{\RtEighteenExcPlainCiHi}{4.898}
\newcommand{\RtEighteenExcPlainCiLo}{4.858}
\newcommand{\RtEighteenLiteral}{18.20.8}
\newcommand{\RtEighteenOpenssl}{3.0.16}
\newcommand{\RtEighteenProbe}{390.5}
\newcommand{\RtEighteenProbeCiHi}{392.9}
\newcommand{\RtEighteenProbeCiLo}{386.9}
\newcommand{\RtEighteenProbeNoStack}{382.7}
\newcommand{\RtEighteenProbeNoStackCiHi}{388.2}
\newcommand{\RtEighteenProbeNoStackCiLo}{380.4}
\newcommand{\RtEighteenStackCapture}{5.397}
\newcommand{\RtEighteenStackCaptureCiHi}{9.122}
\newcommand{\RtEighteenStackCaptureCiLo}{2.657}
\newcommand{\RtEighteenTimer}{0.110}
\newcommand{\RtEighteenTimerCiHi}{0.115}
\newcommand{\RtEighteenTimerCiLo}{0.103}
\newcommand{\RtEighteenVeight}{10}
\newcommand{\RtEighteenVeightFull}{10.2.154.26}
\newcommand{\RtTwentyDecompC}{431.3}
\newcommand{\RtTwentyDecompCCiHi}{431.8}
\newcommand{\RtTwentyDecompCCiLo}{429.2}
\newcommand{\RtTwentyDecompCPct}{97}
\newcommand{\RtTwentyDecompCVersion}{3.0.19}
\newcommand{\RtTwentyDecompExcPct}{1.1}
\newcommand{\RtTwentyDecompLeftover}{7.292}
\newcommand{\RtTwentyDecompLeftoverCiHi}{11.561}
\newcommand{\RtTwentyDecompLeftoverCiLo}{5.325}
\newcommand{\RtTwentyDecompLeftoverPct}{1.6}
\newcommand{\RtTwentyExcNoStack}{2.157}
\newcommand{\RtTwentyExcNoStackCiHi}{2.185}
\newcommand{\RtTwentyExcNoStackCiLo}{2.134}
\newcommand{\RtTwentyExcNodeInternal}{8.635}
\newcommand{\RtTwentyExcNodeInternalCiHi}{9.035}
\newcommand{\RtTwentyExcNodeInternalCiLo}{8.511}
\newcommand{\RtTwentyExcNodeLike}{4.796}
\newcommand{\RtTwentyExcNodeLikeCiHi}{4.815}
\newcommand{\RtTwentyExcNodeLikeCiLo}{4.780}
\newcommand{\RtTwentyExcPlain}{4.874}
\newcommand{\RtTwentyExcPlainCiHi}{4.901}
\newcommand{\RtTwentyExcPlainCiLo}{4.861}
\newcommand{\RtTwentyFourDecompC}{8.447}
\newcommand{\RtTwentyFourDecompCCiHi}{8.469}
\newcommand{\RtTwentyFourDecompCCiLo}{8.412}
\newcommand{\RtTwentyFourDecompCPct}{43}
\newcommand{\RtTwentyFourDecompCVersion}{3.5.8}
\newcommand{\RtTwentyFourDecompExcPct}{21.2}
\newcommand{\RtTwentyFourDecompLeftover}{6.952}
\newcommand{\RtTwentyFourDecompLeftoverCiHi}{7.103}
\newcommand{\RtTwentyFourDecompLeftoverCiLo}{6.803}
\newcommand{\RtTwentyFourDecompLeftoverPct}{35.6}
\newcommand{\RtTwentyFourExcNoStack}{1.809}
\newcommand{\RtTwentyFourExcNoStackCiHi}{1.819}
\newcommand{\RtTwentyFourExcNoStackCiLo}{1.801}
\newcommand{\RtTwentyFourExcNodeInternal}{9.351}
\newcommand{\RtTwentyFourExcNodeInternalCiHi}{9.389}
\newcommand{\RtTwentyFourExcNodeInternalCiLo}{9.315}
\newcommand{\RtTwentyFourExcNodeLike}{4.136}
\newcommand{\RtTwentyFourExcNodeLikeCiHi}{4.145}
\newcommand{\RtTwentyFourExcNodeLikeCiLo}{4.127}
\newcommand{\RtTwentyFourExcPlain}{4.612}
\newcommand{\RtTwentyFourExcPlainCiHi}{4.627}
\newcommand{\RtTwentyFourExcPlainCiLo}{4.587}
\newcommand{\RtTwentyFourLiteral}{24.21.0}
\newcommand{\RtTwentyFourOpenssl}{3.5.8}
\newcommand{\RtTwentyFourProbe}{19.54}
\newcommand{\RtTwentyFourProbeCiHi}{19.66}
\newcommand{\RtTwentyFourProbeCiLo}{19.41}
\newcommand{\RtTwentyFourProbeNoStack}{15.79}
\newcommand{\RtTwentyFourProbeNoStackCiHi}{15.84}
\newcommand{\RtTwentyFourProbeNoStackCiLo}{15.75}
\newcommand{\RtTwentyFourStackCapture}{3.746}
\newcommand{\RtTwentyFourStackCaptureCiHi}{3.813}
\newcommand{\RtTwentyFourStackCaptureCiLo}{3.623}
\newcommand{\RtTwentyFourTimer}{0.066}
\newcommand{\RtTwentyFourTimerCiHi}{0.068}
\newcommand{\RtTwentyFourTimerCiLo}{0.065}
\newcommand{\RtTwentyFourVeight}{13}
\newcommand{\RtTwentyFourVeightFull}{13.6.233.17}
\newcommand{\RtTwentyLiteral}{20.20.2}
\newcommand{\RtTwentyOpenssl}{3.0.19}
\newcommand{\RtTwentyProbe}{443.4}
\newcommand{\RtTwentyProbeCiHi}{446.3}
\newcommand{\RtTwentyProbeCiLo}{441.1}
\newcommand{\RtTwentyProbeNoStack}{437.3}
\newcommand{\RtTwentyProbeNoStackCiHi}{439.6}
\newcommand{\RtTwentyProbeNoStackCiLo}{435.5}
\newcommand{\RtTwentySixDecompC}{8.447}
\newcommand{\RtTwentySixDecompCCiHi}{8.475}
\newcommand{\RtTwentySixDecompCCiLo}{8.432}
\newcommand{\RtTwentySixDecompCPct}{47}
\newcommand{\RtTwentySixDecompCVersion}{3.5.8}
\newcommand{\RtTwentySixDecompExcPct}{19.0}
\newcommand{\RtTwentySixDecompLeftover}{6.145}
\newcommand{\RtTwentySixDecompLeftoverCiHi}{6.249}
\newcommand{\RtTwentySixDecompLeftoverCiLo}{5.988}
\newcommand{\RtTwentySixDecompLeftoverPct}{34.1}
\newcommand{\RtTwentySixExcNoStack}{1.930}
\newcommand{\RtTwentySixExcNoStackCiHi}{1.943}
\newcommand{\RtTwentySixExcNoStackCiLo}{1.923}
\newcommand{\RtTwentySixExcNodeInternal}{6.411}
\newcommand{\RtTwentySixExcNodeInternalCiHi}{6.692}
\newcommand{\RtTwentySixExcNodeInternalCiLo}{6.330}
\newcommand{\RtTwentySixExcNodeLike}{3.421}
\newcommand{\RtTwentySixExcNodeLikeCiHi}{3.565}
\newcommand{\RtTwentySixExcNodeLikeCiLo}{3.403}
\newcommand{\RtTwentySixExcPlain}{2.522}
\newcommand{\RtTwentySixExcPlainCiHi}{2.617}
\newcommand{\RtTwentySixExcPlainCiLo}{2.514}
\newcommand{\RtTwentySixLiteral}{26.10.0}
\newcommand{\RtTwentySixOpenssl}{3.5.8}
\newcommand{\RtTwentySixProbe}{18.01}
\newcommand{\RtTwentySixProbeCiHi}{18.11}
\newcommand{\RtTwentySixProbeCiLo}{17.90}
\newcommand{\RtTwentySixProbeNoStack}{16.16}
\newcommand{\RtTwentySixProbeNoStackCiHi}{16.25}
\newcommand{\RtTwentySixProbeNoStackCiLo}{16.13}
\newcommand{\RtTwentySixStackCapture}{1.767}
\newcommand{\RtTwentySixStackCaptureCiHi}{1.997}
\newcommand{\RtTwentySixStackCaptureCiLo}{1.724}
\newcommand{\RtTwentySixTimer}{0.067}
\newcommand{\RtTwentySixTimerCiHi}{0.068}
\newcommand{\RtTwentySixTimerCiLo}{0.066}
\newcommand{\RtTwentySixVeight}{14}
\newcommand{\RtTwentySixVeightFull}{14.6.202.34}
\newcommand{\RtTwentyStackCapture}{4.917}
\newcommand{\RtTwentyStackCaptureCiHi}{8.704}
\newcommand{\RtTwentyStackCaptureCiLo}{3.370}
\newcommand{\RtTwentyThreeDecompC}{394.4}
\newcommand{\RtTwentyThreeDecompCCiHi}{395.3}
\newcommand{\RtTwentyThreeDecompCCiLo}{389.2}
\newcommand{\RtTwentyThreeDecompCPct}{103}
\newcommand{\RtTwentyThreeDecompCVersion}{3.0.16}
\newcommand{\RtTwentyThreeDecompExcPct}{1.0}
\newcommand{\RtTwentyThreeDecompLeftover}{-14.88}
\newcommand{\RtTwentyThreeDecompLeftoverCiHi}{-8.66}
\newcommand{\RtTwentyThreeDecompLeftoverCiLo}{-18.64}
\newcommand{\RtTwentyThreeDecompLeftoverPct}{-3.9}
\newcommand{\RtTwentyThreeExcNoStack}{1.682}
\newcommand{\RtTwentyThreeExcNoStackCiHi}{1.683}
\newcommand{\RtTwentyThreeExcNoStackCiLo}{1.674}
\newcommand{\RtTwentyThreeExcNodeInternal}{7.718}
\newcommand{\RtTwentyThreeExcNodeInternalCiHi}{7.826}
\newcommand{\RtTwentyThreeExcNodeInternalCiLo}{7.662}
\newcommand{\RtTwentyThreeExcNodeLike}{3.959}
\newcommand{\RtTwentyThreeExcNodeLikeCiHi}{3.973}
\newcommand{\RtTwentyThreeExcNodeLikeCiLo}{3.951}
\newcommand{\RtTwentyThreeExcPlain}{4.526}
\newcommand{\RtTwentyThreeExcPlainCiHi}{4.562}
\newcommand{\RtTwentyThreeExcPlainCiLo}{4.516}
\newcommand{\RtTwentyThreeLiteral}{23.11.1}
\newcommand{\RtTwentyThreeOpenssl}{3.0.16}
\newcommand{\RtTwentyThreeProbe}{383.5}
\newcommand{\RtTwentyThreeProbeCiHi}{386.8}
\newcommand{\RtTwentyThreeProbeCiLo}{379.9}
\newcommand{\RtTwentyThreeProbeNoStack}{379.9}
\newcommand{\RtTwentyThreeProbeNoStackCiHi}{384.9}
\newcommand{\RtTwentyThreeProbeNoStackCiLo}{375.1}
\newcommand{\RtTwentyThreeStackCapture}{4.675}
\newcommand{\RtTwentyThreeStackCaptureCiHi}{9.334}
\newcommand{\RtTwentyThreeStackCaptureCiLo}{-1.535}
\newcommand{\RtTwentyThreeTimer}{0.068}
\newcommand{\RtTwentyThreeTimerCiHi}{0.069}
\newcommand{\RtTwentyThreeTimerCiLo}{0.067}
\newcommand{\RtTwentyThreeVeight}{12}
\newcommand{\RtTwentyThreeVeightFull}{12.9.202.28}
\newcommand{\RtTwentyTimer}{0.070}
\newcommand{\RtTwentyTimerCiHi}{0.083}
\newcommand{\RtTwentyTimerCiLo}{0.068}
\newcommand{\RtTwentyTwoDecompC}{8.447}
\newcommand{\RtTwentyTwoDecompCCiHi}{8.475}
\newcommand{\RtTwentyTwoDecompCCiLo}{8.412}
\newcommand{\RtTwentyTwoDecompCPct}{44}
\newcommand{\RtTwentyTwoDecompCVersion}{3.5.8}
\newcommand{\RtTwentyTwoDecompExcPct}{21.2}
\newcommand{\RtTwentyTwoDecompLeftover}{6.515}
\newcommand{\RtTwentyTwoDecompLeftoverCiHi}{6.678}
\newcommand{\RtTwentyTwoDecompLeftoverCiLo}{6.427}
\newcommand{\RtTwentyTwoDecompLeftoverPct}{34.3}
\newcommand{\RtTwentyTwoExcNoStack}{1.710}
\newcommand{\RtTwentyTwoExcNoStackCiHi}{1.719}
\newcommand{\RtTwentyTwoExcNoStackCiLo}{1.708}
\newcommand{\RtTwentyTwoExcNodeInternal}{7.735}
\newcommand{\RtTwentyTwoExcNodeInternalCiHi}{7.889}
\newcommand{\RtTwentyTwoExcNodeInternalCiLo}{7.641}
\newcommand{\RtTwentyTwoExcNodeLike}{4.021}
\newcommand{\RtTwentyTwoExcNodeLikeCiHi}{4.044}
\newcommand{\RtTwentyTwoExcNodeLikeCiLo}{3.999}
\newcommand{\RtTwentyTwoExcPlain}{4.553}
\newcommand{\RtTwentyTwoExcPlainCiHi}{4.569}
\newcommand{\RtTwentyTwoExcPlainCiLo}{4.538}
\newcommand{\RtTwentyTwoLiteral}{22.23.3}
\newcommand{\RtTwentyTwoOpenssl}{3.5.8}
\newcommand{\RtTwentyTwoProbe}{18.98}
\newcommand{\RtTwentyTwoProbeCiHi}{19.11}
\newcommand{\RtTwentyTwoProbeCiLo}{18.91}
\newcommand{\RtTwentyTwoProbeNoStack}{15.18}
\newcommand{\RtTwentyTwoProbeNoStackCiHi}{15.27}
\newcommand{\RtTwentyTwoProbeNoStackCiLo}{15.15}
\newcommand{\RtTwentyTwoStackCapture}{3.805}
\newcommand{\RtTwentyTwoStackCaptureCiHi}{3.946}
\newcommand{\RtTwentyTwoStackCaptureCiLo}{3.645}
\newcommand{\RtTwentyTwoTimer}{0.061}
\newcommand{\RtTwentyTwoTimerCiHi}{0.063}
\newcommand{\RtTwentyTwoTimerCiLo}{0.061}
\newcommand{\RtTwentyTwoVeight}{12}
\newcommand{\RtTwentyTwoVeightFull}{12.4.254.21}
\newcommand{\RtTwentyVeight}{11}
\newcommand{\RtTwentyVeightFull}{11.3.244.8}
\newcommand{\XhCallsPerInvocation}{40000}
\newcommand{\XhCreateSecretKey}{1.261}
\newcommand{\XhCreateSecretKeyCiHi}{1.270}
\newcommand{\XhCreateSecretKeyCiLo}{1.253}
\newcommand{\XhCvMax}{20.9}
\newcommand{\XhDecodeOnly}{1.795}
\newcommand{\XhDecodeOnlyCiHi}{1.803}
\newcommand{\XhDecodeOnlyCiLo}{1.787}
\newcommand{\XhDecompExplainedPct}{75.8}
\newcommand{\XhHmacKeyObject}{2.974}
\newcommand{\XhHmacKeyObjectCiHi}{2.986}
\newcommand{\XhHmacKeyObjectCiLo}{2.961}
\newcommand{\XhHmacString}{3.105}
\newcommand{\XhHmacStringCiHi}{3.113}
\newcommand{\XhHmacStringCiLo}{3.086}
\newcommand{\XhInvocations}{15}
\newcommand{\XhJwtPreparsed}{9.083}
\newcommand{\XhJwtPreparsedCiHi}{9.118}
\newcommand{\XhJwtPreparsedCiLo}{9.050}
\newcommand{\XhJwtString}{36.48}
\newcommand{\XhJwtStringCiHi}{36.75}
\newcommand{\XhJwtStringCiLo}{36.32}
\newcommand{\XhJwtStringSafe}{11.09}
\newcommand{\XhJwtStringSafeCiHi}{11.33}
\newcommand{\XhJwtStringSafeCiLo}{11.04}
\newcommand{\XhNode}{22.22.2}
\newcommand{\XhOpenssl}{3.5.5}
\newcommand{\XhProbeShareOfCall}{53}
\newcommand{\XhProbeSucceeds}{13.66}
\newcommand{\XhProbeSucceedsCiHi}{13.69}
\newcommand{\XhProbeSucceedsCiLo}{13.65}
\newcommand{\XhProbeThrows}{19.42}
\newcommand{\XhProbeThrowsCiHi}{19.46}
\newcommand{\XhProbeThrowsCiLo}{19.35}
\newcommand{\XhStringPenalty}{27.30}
\newcommand{\XhStringPenaltyCiHi}{27.61}
\newcommand{\XhStringPenaltyCiLo}{27.13}
\newcommand{\XhTimerOverhead}{0.058}
\newcommand{\XhTimerOverheadCiHi}{0.061}
\newcommand{\XhTimerOverheadCiLo}{0.058}
\newcommand{\XlibCallsPerInvocation}{20000}
\newcommand{\XlibCvMax}{20.2}
\newcommand{\XlibCvMedian}{2.8}
\newcommand{\XlibFastjwtNewSslBackend}{OpenSSL 3.5.8}
\newcommand{\XlibFastjwtNewSslPenalty}{-0.068}
\newcommand{\XlibFastjwtNewSslPenaltyCiHi}{0.197}
\newcommand{\XlibFastjwtNewSslPenaltyCiLo}{-0.134}
\newcommand{\XlibFastjwtNewSslPenaltyExcludesZero}{no}
\newcommand{\XlibFastjwtNewSslPreparsed}{7.934}
\newcommand{\XlibFastjwtNewSslPreparsedCiHi}{7.992}
\newcommand{\XlibFastjwtNewSslPreparsedCiLo}{7.866}
\newcommand{\XlibFastjwtNewSslRatio}{1.00}
\newcommand{\XlibFastjwtNewSslRuntime}{node v26.10.0}
\newcommand{\XlibFastjwtNewSslString}{7.917}
\newcommand{\XlibFastjwtNewSslStringCiHi}{8.044}
\newcommand{\XlibFastjwtNewSslStringCiLo}{7.851}
\newcommand{\XlibFastjwtNewSslVersion}{6.3.3}
\newcommand{\XlibFastjwtOldSslBackend}{OpenSSL 3.0.19}
\newcommand{\XlibFastjwtOldSslPenalty}{0.070}
\newcommand{\XlibFastjwtOldSslPenaltyCiHi}{0.111}
\newcommand{\XlibFastjwtOldSslPenaltyCiLo}{-0.116}
\newcommand{\XlibFastjwtOldSslPenaltyExcludesZero}{no}
\newcommand{\XlibFastjwtOldSslPreparsed}{7.570}
\newcommand{\XlibFastjwtOldSslPreparsedCiHi}{7.642}
\newcommand{\XlibFastjwtOldSslPreparsedCiLo}{7.445}
\newcommand{\XlibFastjwtOldSslRatio}{1.00}
\newcommand{\XlibFastjwtOldSslRuntime}{node v20.20.2}
\newcommand{\XlibFastjwtOldSslString}{7.564}
\newcommand{\XlibFastjwtOldSslStringCiHi}{7.666}
\newcommand{\XlibFastjwtOldSslStringCiLo}{7.487}
\newcommand{\XlibFastjwtOldSslVersion}{6.3.3}
\newcommand{\XlibGolangjwtBackend}{Go crypto/hmac}
\newcommand{\XlibGolangjwtPenalty}{0.034}
\newcommand{\XlibGolangjwtPenaltyCiHi}{0.117}
\newcommand{\XlibGolangjwtPenaltyCiLo}{-0.042}
\newcommand{\XlibGolangjwtPenaltyExcludesZero}{no}
\newcommand{\XlibGolangjwtPreparsed}{4.764}
\newcommand{\XlibGolangjwtPreparsedCiHi}{4.832}
\newcommand{\XlibGolangjwtPreparsedCiLo}{4.719}
\newcommand{\XlibGolangjwtRatio}{1.01}
\newcommand{\XlibGolangjwtRuntime}{go1.24.7}
\newcommand{\XlibGolangjwtString}{4.810}
\newcommand{\XlibGolangjwtStringCiHi}{4.881}
\newcommand{\XlibGolangjwtStringCiLo}{4.734}
\newcommand{\XlibGolangjwtVersion}{5.3.1}
\newcommand{\XlibInvocations}{15}
\newcommand{\XlibJjwtBackend}{JCA SunJCE}
\newcommand{\XlibJjwtPenalty}{-0.406}
\newcommand{\XlibJjwtPenaltyCiHi}{0.648}
\newcommand{\XlibJjwtPenaltyCiLo}{-0.715}
\newcommand{\XlibJjwtPenaltyExcludesZero}{no}
\newcommand{\XlibJjwtPreparsed}{12.33}
\newcommand{\XlibJjwtPreparsedCiHi}{12.78}
\newcommand{\XlibJjwtPreparsedCiLo}{12.04}
\newcommand{\XlibJjwtRatio}{0.99}
\newcommand{\XlibJjwtRuntime}{java 21.0.10}
\newcommand{\XlibJjwtString}{12.15}
\newcommand{\XlibJjwtStringCiHi}{13.12}
\newcommand{\XlibJjwtStringCiLo}{11.98}
\newcommand{\XlibJjwtVersion}{0.13.0}
\newcommand{\XlibJoseNewSslBackend}{OpenSSL 3.5.8}
\newcommand{\XlibJoseNewSslPenalty}{24.80}
\newcommand{\XlibJoseNewSslPenaltyCiHi}{26.51}
\newcommand{\XlibJoseNewSslPenaltyCiLo}{21.91}
\newcommand{\XlibJoseNewSslPenaltyExcludesZero}{yes}
\newcommand{\XlibJoseNewSslPreparsed}{73.98}
\newcommand{\XlibJoseNewSslPreparsedCiHi}{75.81}
\newcommand{\XlibJoseNewSslPreparsedCiLo}{71.27}
\newcommand{\XlibJoseNewSslRatio}{1.34}
\newcommand{\XlibJoseNewSslRuntime}{node v26.10.0}
\newcommand{\XlibJoseNewSslString}{98.89}
\newcommand{\XlibJoseNewSslStringCiHi}{100.02}
\newcommand{\XlibJoseNewSslStringCiLo}{95.81}
\newcommand{\XlibJoseNewSslVersion}{6.2.12}
\newcommand{\XlibJoseOldSslBackend}{OpenSSL 3.0.19}
\newcommand{\XlibJoseOldSslPenalty}{42.66}
\newcommand{\XlibJoseOldSslPenaltyCiHi}{45.89}
\newcommand{\XlibJoseOldSslPenaltyCiLo}{38.50}
\newcommand{\XlibJoseOldSslPenaltyExcludesZero}{yes}
\newcommand{\XlibJoseOldSslPreparsed}{78.42}
\newcommand{\XlibJoseOldSslPreparsedCiHi}{81.65}
\newcommand{\XlibJoseOldSslPreparsedCiLo}{75.86}
\newcommand{\XlibJoseOldSslRatio}{1.55}
\newcommand{\XlibJoseOldSslRuntime}{node v20.20.2}
\newcommand{\XlibJoseOldSslString}{121.2}
\newcommand{\XlibJoseOldSslStringCiHi}{124.5}
\newcommand{\XlibJoseOldSslStringCiLo}{119.9}
\newcommand{\XlibJoseOldSslVersion}{6.2.12}
\newcommand{\XlibJsonwebtokenNewSslBackend}{OpenSSL 3.5.8}
\newcommand{\XlibJsonwebtokenNewSslPenalty}{24.56}
\newcommand{\XlibJsonwebtokenNewSslPenaltyCiHi}{24.82}
\newcommand{\XlibJsonwebtokenNewSslPenaltyCiLo}{24.07}
\newcommand{\XlibJsonwebtokenNewSslPenaltyExcludesZero}{yes}
\newcommand{\XlibJsonwebtokenNewSslPreparsed}{9.844}
\newcommand{\XlibJsonwebtokenNewSslPreparsedCiHi}{9.991}
\newcommand{\XlibJsonwebtokenNewSslPreparsedCiLo}{9.634}
\newcommand{\XlibJsonwebtokenNewSslRatio}{3.48}
\newcommand{\XlibJsonwebtokenNewSslRuntime}{node v26.10.0}
\newcommand{\XlibJsonwebtokenNewSslString}{34.22}
\newcommand{\XlibJsonwebtokenNewSslStringCiHi}{34.70}
\newcommand{\XlibJsonwebtokenNewSslStringCiLo}{33.98}
\newcommand{\XlibJsonwebtokenNewSslVersion}{9.0.3}
\newcommand{\XlibJsonwebtokenOldSslBackend}{OpenSSL 3.0.19}
\newcommand{\XlibJsonwebtokenOldSslPenalty}{465.9}
\newcommand{\XlibJsonwebtokenOldSslPenaltyCiHi}{467.8}
\newcommand{\XlibJsonwebtokenOldSslPenaltyCiLo}{465.2}
\newcommand{\XlibJsonwebtokenOldSslPenaltyExcludesZero}{yes}
\newcommand{\XlibJsonwebtokenOldSslPreparsed}{10.23}
\newcommand{\XlibJsonwebtokenOldSslPreparsedCiHi}{10.45}
\newcommand{\XlibJsonwebtokenOldSslPreparsedCiLo}{10.16}
\newcommand{\XlibJsonwebtokenOldSslRatio}{46.55}
\newcommand{\XlibJsonwebtokenOldSslRuntime}{node v20.20.2}
\newcommand{\XlibJsonwebtokenOldSslString}{476.2}
\newcommand{\XlibJsonwebtokenOldSslStringCiHi}{477.5}
\newcommand{\XlibJsonwebtokenOldSslStringCiLo}{475.4}
\newcommand{\XlibJsonwebtokenOldSslVersion}{9.0.3}
\newcommand{\XlibJwtOverJoseOldSsl}{11}
\newcommand{\XlibLibCount}{seven}
\newcommand{\XlibNimbusBackend}{JCA SunJCE}
\newcommand{\XlibNimbusPenalty}{0.304}
\newcommand{\XlibNimbusPenaltyCiHi}{1.016}
\newcommand{\XlibNimbusPenaltyCiLo}{-0.098}
\newcommand{\XlibNimbusPenaltyExcludesZero}{no}
\newcommand{\XlibNimbusPreparsed}{4.229}
\newcommand{\XlibNimbusPreparsedCiHi}{4.418}
\newcommand{\XlibNimbusPreparsedCiLo}{4.134}
\newcommand{\XlibNimbusRatio}{1.03}
\newcommand{\XlibNimbusRuntime}{java 21.0.10}
\newcommand{\XlibNimbusString}{4.358}
\newcommand{\XlibNimbusStringCiHi}{5.799}
\newcommand{\XlibNimbusStringCiLo}{4.247}
\newcommand{\XlibNimbusVersion}{10.10}
\newcommand{\XlibPenaltyLibCount}{two}
\newcommand{\XlibPyjwtBackend}{hashlib, OpenSSL 3.0.13}
\newcommand{\XlibPyjwtPenalty}{-0.103}
\newcommand{\XlibPyjwtPenaltyCiHi}{0.338}
\newcommand{\XlibPyjwtPenaltyCiLo}{-0.457}
\newcommand{\XlibPyjwtPenaltyExcludesZero}{no}
\newcommand{\XlibPyjwtPreparsed}{41.28}
\newcommand{\XlibPyjwtPreparsedCiHi}{41.49}
\newcommand{\XlibPyjwtPreparsedCiLo}{41.08}
\newcommand{\XlibPyjwtRatio}{1.00}
\newcommand{\XlibPyjwtRuntime}{cpython 3.11.15}
\newcommand{\XlibPyjwtString}{41.33}
\newcommand{\XlibPyjwtStringCiHi}{41.40}
\newcommand{\XlibPyjwtStringCiLo}{41.22}
\newcommand{\XlibPyjwtVersion}{2.15.0}
\newcommand{\XlibZeroPenaltyLibCount}{five}
\newcommand{\XmCallsPerInvocation}{20000}
\newcommand{\XmEighteenArch}{x64}
\newcommand{\XmEighteenDecodeOnly}{2.32}
\newcommand{\XmEighteenHmacString}{4.64}
\newcommand{\XmEighteenJwtPreparsed}{12.66}
\newcommand{\XmEighteenJwtPreparsedCiHi}{12.81}
\newcommand{\XmEighteenJwtPreparsedCiLo}{12.57}
\newcommand{\XmEighteenJwtString}{694.5}
\newcommand{\XmEighteenJwtStringCiHi}{696.9}
\newcommand{\XmEighteenJwtStringCiLo}{690.4}
\newcommand{\XmEighteenOpenssl}{3.0.16}
\newcommand{\XmEighteenProbeSucceeds}{181.25}
\newcommand{\XmEighteenProbeThrows}{607.1}
\newcommand{\XmEighteenProbeThrowsCiHi}{613.7}
\newcommand{\XmEighteenProbeThrowsCiLo}{602.8}
\newcommand{\XmInvocations}{10}
\newcommand{\XmOrderingMatchesArm}{a different}
\newcommand{\XmSeriesSplitHolds}{holds}
\newcommand{\XmSpanDecodeOnly}{1.36}
\newcommand{\XmSpanHmacString}{1.21}
\newcommand{\XmSpanJwtPreparsed}{1.31}
\newcommand{\XmSpanProbeSucceeds}{14.75}
\newcommand{\XmSpanProbeThrows}{30.50}
\newcommand{\XmTwentyArch}{x64}
\newcommand{\XmTwentyDecodeOnly}{2.32}
\newcommand{\XmTwentyHmacString}{4.15}
\newcommand{\XmTwentyJwtPreparsed}{11.89}
\newcommand{\XmTwentyJwtPreparsedCiHi}{11.99}
\newcommand{\XmTwentyJwtPreparsedCiLo}{11.68}
\newcommand{\XmTwentyJwtString}{632.0}
\newcommand{\XmTwentyJwtStringCiHi}{633.7}
\newcommand{\XmTwentyJwtStringCiLo}{628.1}
\newcommand{\XmTwentyOpenssl}{3.0.19}
\newcommand{\XmTwentyProbeSucceeds}{196.98}
\newcommand{\XmTwentyProbeThrows}{556.3}
\newcommand{\XmTwentyProbeThrowsCiHi}{559.6}
\newcommand{\XmTwentyProbeThrowsCiLo}{552.2}
\newcommand{\XmTwentySixGlibcArch}{x64}
\newcommand{\XmTwentySixGlibcDecodeOnly}{1.71}
\newcommand{\XmTwentySixGlibcHmacString}{3.84}
\newcommand{\XmTwentySixGlibcJwtPreparsed}{9.701}
\newcommand{\XmTwentySixGlibcJwtPreparsedCiHi}{9.790}
\newcommand{\XmTwentySixGlibcJwtPreparsedCiLo}{9.528}
\newcommand{\XmTwentySixGlibcJwtString}{39.64}
\newcommand{\XmTwentySixGlibcJwtStringCiHi}{39.93}
\newcommand{\XmTwentySixGlibcJwtStringCiLo}{39.21}
\newcommand{\XmTwentySixGlibcOpenssl}{3.5.7}
\newcommand{\XmTwentySixGlibcProbeSucceeds}{13.35}
\newcommand{\XmTwentySixGlibcProbeThrows}{19.91}
\newcommand{\XmTwentySixGlibcProbeThrowsCiHi}{20.06}
\newcommand{\XmTwentySixGlibcProbeThrowsCiLo}{19.75}
\newcommand{\XmTwentySixMuslArch}{x64}
\newcommand{\XmTwentySixMuslDecodeOnly}{1.82}
\newcommand{\XmTwentySixMuslHmacString}{4.57}
\newcommand{\XmTwentySixMuslJwtPreparsed}{10.94}
\newcommand{\XmTwentySixMuslJwtPreparsedCiHi}{11.12}
\newcommand{\XmTwentySixMuslJwtPreparsedCiLo}{10.67}
\newcommand{\XmTwentySixMuslJwtString}{74.25}
\newcommand{\XmTwentySixMuslJwtStringCiHi}{74.95}
\newcommand{\XmTwentySixMuslJwtStringCiLo}{73.55}
\newcommand{\XmTwentySixMuslOpenssl}{3.5.7}
\newcommand{\XmTwentySixMuslProbeSucceeds}{18.86}
\newcommand{\XmTwentySixMuslProbeThrows}{42.13}
\newcommand{\XmTwentySixMuslProbeThrowsCiHi}{42.70}
\newcommand{\XmTwentySixMuslProbeThrowsCiLo}{41.90}

\newcommand{\LibraryName}{\texttt{jsonwebtoken}}
\newcommand{\LibraryTag}{tag~v9.0.3}
\renewcommand{\thesection}{S\arabic{section}}
\renewcommand{\thetable}{S\arabic{table}}
\renewcommand{\thefigure}{S\arabic{figure}}
\emergencystretch=3em
\journal{Journal of Systems and Software}

\begin{document}
\begin{frontmatter}
\title{Supplementary material for ``It Is Not the Signature: An Empirical Study
of a Hidden Key-Parse Cost in JSON Web Token Validation''}
\author{Jannatul Ferdous Katha, Tasmia Tahmid Prova, Md Kishor Morol, Tze Hui
Liew, Dip Nandi}
\begin{abstract}
This supplement contains material moved from the main paper for length: the
extended related work, measurement-protocol details, additional analyses for
RQ1 to RQ3, the full discussion of threats to validity, and a design note on
relocating validation to the service mesh. Section, table and figure numbers
without an ``S'' prefix refer to the main paper. Every value is generated
from the replication package, as in the main paper.
\end{abstract}
\end{frontmatter}

%% ======================================================================
\section{Extended related work}

This section gives the full discussion of related work that Section~\ref{sec:related}
of the main paper summarises.

\paragraph{Where the cost is usually attributed} Published measurement of
layered enforcement in microservices concentrates on transport and proxy
layers. \citet{zhu2023dissecting} decompose service-mesh sidecar
overhead and show that the dominant contributors depend on configuration and
workload rather than being a fixed tax; we reach an analogous conclusion one
layer up, for application-level validation. Application-level token validation
is itself thinly evidenced, and is commonly reported as a difference between
scenario averages, which attributes to validation every difference between two
runs. We were unable to locate a study that measures the key-resolution step
directly;
we state that as an absence we searched for rather than as a claim about the
whole literature.

\paragraph{Cryptographic cost has been over-attributed before} \citet{coarfa2006performance} remains the clearest precedent: measuring TLS
web servers rather than primitives, they found the public-key operation was a
minority of server cost. Work on cryptographic library performance since has
largely stayed at the level of the primitive or the protocol
exchange~\citep{naser2020performance,sikeridis2020assessing}, with \citet{tuveri2019engines} a notable exception in examining how a library
loads and dispatches implementations rather than how fast those
implementations run. Our result concerns neither: it is the cost of getting key
material \emph{into} the library.

\paragraph{Performance bugs as a class} \citet{jin2012understanding} studied real-world performance bugs across five
code bases and found that most follow a small number of recurring patterns (among them uncoordinated and skippable function calls) and that fixes are
usually local. \citet{zaman2012qualitative} found, from 400
Firefox and Chrome reports, that performance bugs differ from other bugs in
how they are reported, discussed and fixed, in part because reproducing them
is hard; \citet{nistor2013discovering} report that most are found
by code reasoning rather than by observing symptoms. \citet{liu2014characterizing} reach similar conclusions for mobile
applications. Detection work targets structural symptoms: redundant traversals
with similar memory-access patterns~\citep{nistor2013toddler}, or predicates
that separate slow from fast runs~\citep{song2014statistical}. The defect here
is of the ``skippable call'' kind Jin et al.\ describe, but it has none of the
structural signatures those detectors look for: it is one call, executed once
per request, whose failure is its only observable effect. It is also the shape
that makes a performance bug hard to report: our own upstream report was
preceded by a two-year-old report of the symptom with no diagnosed cause
(Section~\ref{sec:upstream}).

\paragraph{JavaScript and Node.js performance} \citet{selakovic2016performance} classified performance issues in
JavaScript projects and found that inefficient API usage (correct calls made
in a costly way) is among the most common root causes, and that the fixes
are typically a few lines. Their later profiler~\citep{selakovic2017actionable}
targets precisely the \emph{order} in which a program evaluates conditions, the
axis on which this defect lives: a check that almost always fails is evaluated
before one that almost always succeeds. \citet{gong2015jitprof}
locate code that defeats JIT optimisation; exception-heavy control flow is one
such pattern, and Section~\ref{sec:granularity} measures what it costs here.
Work on the Node.js event loop shows why a per-request synchronous cost matters
beyond its magnitude: a single slow handler delays every other
client~\citep{davis2018sense,staicu2018freezing}. We make no availability or
capacity claim, but the mechanism by which a synchronous per-request cost
propagates is the same.

\paragraph{Version and environment variation} That software performance
changes across versions in ways maintainers do not anticipate is established:
\citet{kaltenecker2023performance} track performance across
configurable systems' release histories, and \citet{chen2017exploratory} study the code changes that introduce
regressions. \citet{laaber2018evaluation} assess whether
microbenchmark suites detect such changes at all, and \citet{stefan2017unit} find performance unit testing rare in practice.
Container images make this concrete by pinning a dependency set implicitly:
\citet{cito2017empirical} characterise how the ecosystem uses them,
\citet{zerouali2019outdated} show that images routinely ship
outdated packages, and \citet{malka2026docker} show that an image
tag does not guarantee a reproducible build. For OpenSSL~3 specifically, the
maintainers of a major downstream library have stated publicly that parsing and
key loading regressed substantially against 1.1.1~\citep{pyca2026stateofopenssl};
that statement is not peer reviewed and was not produced under a protocol we
can assess, and Section~\ref{sec:rq2-causal} supplies a controlled measurement
for the one path this paper concerns. Section~\ref{sec:rq2} is an instance with
a measured cost attached.

\paragraph{Measurement protocol is the thing at stake} The methodological core
of this paper is not new. \citet{georges2007statistically}
established that a single execution of a managed-runtime benchmark is not a
measurement, and that intervals must be computed across executions rather than
within one; \citet{kalibera2013rigorous} show how to spend a
fixed budget across levels of repetition, which is the reasoning behind our
choice of many short processes over few long ones. \citet{mytkowicz2009wrong} showed that environmental details unrelated to
the code under study can produce effects larger than the one being measured,
without any obvious mistake, and \citet{curtsinger2013stabilizer} that memory layout alone can do so.
\citet{barrett2017warmup} found that virtual machines frequently
do not reach a steady state, \citet{traini2023steady} that
developers misjudge when warmup ends, and \citet{costa2021wrong}
catalogue how microbenchmarks go wrong in practice. \citet{mytkowicz2010profilers} showed that sampling profilers disagree on
which method is hot, which is why Section~\ref{sec:rq1} uses the profiler only
as corroboration of a result obtained by differencing and relies on ordering
rather than on its percentages. Our contribution is not to these methods but is
an instance of what they warn about. The attribution this paper overturns has
exactly the structure Mytkowicz et al.\ describe: a second variable, the
bundled OpenSSL version, moving unrecorded alongside the one under study, at
a granularity coarse enough to hide it (Section~\ref{sec:rq2}).

\paragraph{Environment as a confounder} \citet{laaber2019cloud}
quantify how far cloud environments degrade microbenchmark reliability, and
\citet{leitner2016patterns} characterise performance variation
across public infrastructure providers. Their concern is variance between
nominally identical environments. Ours is adjacent and, we think, less
examined: a \emph{systematic} difference between environments that are not
identical and are routinely treated as though they were: two container base
images differing in a bundled library version. Laaber et al.\ are also the
reason the second-architecture repetition in Section~\ref{sec:rq2-arch} is
reported as an ordering check rather than as a replication of magnitudes.

\paragraph{Cryptographic interfaces are hard to use correctly} The defect
reported here is not a cryptographic weakness but a consequence of how a
cryptographic interface is invoked, which is a well-documented class. \citet{lazar2014why} found that most of the cryptographic vulnerabilities
they studied were misuses of libraries by applications rather than bugs in the
libraries. \citet{egele2013empirical} found cryptographic misuse
widespread in deployed applications and \citet{georgiev2012dangerous} found certificate validation broken across a
wide range of non-browser software; \citet{nadi2016jumping}
attribute much of it to the interfaces themselves rather than to developer
inattention, a position \citet{green2016developers} put directly
and \citet{indela2016helping} respond to with interface proposals;
and \citet{acar2017comparing} show experimentally that API design
materially changes whether developers produce correct code. Detection has been
automated by specifying correct usage, as CrySL does~\citep{kruger2018crysl},
and by static analysis at scale~\citep{rahaman2019cryptoguard}. Our finding sits
within that line with one difference: the interface here admits a use that is
\emph{correct}, and costly only for a reason invisible at the call site. A
CrySL-style rule could express ``pass a key object, not a string'', but the
rule sets we know of encode security properties only, so none would report it.

\paragraph{JWT and JOSE libraries} Security analyses of JWT implementations
consistently find key-handling defects. \citet{detering2017jose} demonstrated attacks, key confusion among them,
against several popular JOSE libraries; \citet{xu2023jwtkey}
automated the detection of cryptographic vulnerabilities in applications that
use JWT; \citet{yang2026tokentimebomb} evaluated 43 JWT libraries
across 16 programming languages for vulnerabilities; and \citet{shatnawi2024python} surveyed Python JWT libraries for security
warnings and adoption. RFC 8725~\citep{rfc8725} codifies the resulting practice.
All of this work treats a library's handling of key material as a security
surface. Section~\ref{sec:rq4} treats the same surface as a cost site, across
seven libraries, and Section~\ref{sec:fix} shows the two views are coupled: the
control flow that costs time in one library is the control flow that enforces
its key-type check.

\paragraph{Mining public repositories} Section~\ref{sec:prevalence} draws on
code hosted publicly, and the hazards of doing so are catalogued by
\citet{kalliamvakou2014promises}, among them that many
repositories are not software projects, that activity is not uniformly
distributed, and that a search result is not a random sample. Our sampling
frame is stated in Section~\ref{sec:frames} with those hazards in view, which
is why no population proportion is computed from it.

\paragraph{The cost of exceptions} That exceptions are expensive in managed
runtimes, and that their cost is concentrated in construction rather than in
unwinding, is long established: \citet{ogasawara2006edo}
optimise around exception paths in a Java virtual machine, and \citet{renwick2019lowcost} quantify the cost in a compiled setting.
\citet{parravicini2021cost} provide the closest
available cost-attribution baseline for the V8 engine, though on speculation
rather than exceptions. Empirical studies of how exceptions are \emph{used}
find handlers that discard what they catch to be common: \citet{cabral2007exception} report that handlers in the Java and .NET
code they studied rarely attempt recovery, \citet{kery2016examining} find that handlers mostly log, print, return or
rethrow, and \citet{depadua2017prevalence} measure the
prevalence of handling anti-patterns, as \citet{weimer2004finding} did earlier for error-handling mistakes. That
literature treats a swallowed exception as a correctness and maintainability
smell. The handler at the centre of this paper is a swallowed exception used as
control flow on a hot path, and it is a performance defect rather than a
correctness one: the fallback it guards is correct. None of these studies
measures such a handler's cost; Section~\ref{sec:granularity} does, and finds
the exception itself the smaller part.

\paragraph{Key-type confusion} That a token declaring an HMAC algorithm must
not be validated against asymmetric key material is long established; RFC
8725~\citep{rfc8725} describes the attack and its mitigations. The library
studied here has been hardened against this class twice. Versions before 4.2.2
accepted an HMAC token verified with an RSA or EC public key used as the
shared secret (CVE-2015-9235)~\citep{ghsa2015jsonwebtoken}, and versions up to
8.5.1 could still be led to verify forged tokens from RSA to HMAC through
their key-retrieval path (CVE-2022-23541), alongside two related defects in
key-type and default-algorithm handling (CVE-2022-23539,
CVE-2022-23540)~\citep{auth0bulletin2022}; all three were fixed in
version~9.0.0. Section~\ref{sec:security} reports no new attack. It reports
that the second hardening introduced the control flow whose cost this paper
measures, that the defence depends on that control flow, and that removing the
cost naively removes the defence.

\subsection*{What we searched for and did not find}

Three areas bear directly on this result and we were unable to locate
verifiable peer-reviewed work in any of them. We state each as an absence we
searched for rather than as a claim about the whole literature; each is also,
in part, why the misattribution corrected here survived as long as it did.
Two adjacent questions (the cost of error construction in V8, and where
key-parsing cost lies inside OpenSSL) were equally unstudied, and this paper
measures both directly (Sections~\ref{sec:granularity}
and~\ref{sec:rq2-causal}).

\emph{OpenSSL 3.x performance regressions.} We found no peer-reviewed study of
performance regressions across the OpenSSL 3.x series. Evidence exists as a
maintainer statement~\citep{pyca2026stateofopenssl}, as project issues (for
instance \texttt{openssl/openssl} issues \#17627 and \#20698, which concern
performance of the 3.x provider and decoder machinery) and as mailing-list
discussion. None was produced under a measurement protocol we can assess, so
Section~\ref{sec:rq2-causal} measures the path directly rather than adopting
the mechanism those sources propose.

\emph{Base image to bundled library to application performance.} We found no
academic study tracing that chain. \citet{malka2026docker} address
build reproducibility, \citet{zerouali2019outdated} package
staleness and \citet{cito2017empirical} ecosystem practice; none
measures how the library version an image pins changes application
performance.

\emph{Cryptographic API misuse as a performance defect.} The misuse literature
cited above treats misuse exclusively as a security problem. We found no work
treating a correct-but-costly use of a cryptographic interface as a latency or
throughput defect, which is precisely what this paper reports.

%% ======================================================================
\section{Measurement protocol details}

This section gives the details of the measurement protocol summarised in
Section~\ref{sec:unit}, the full descriptions of the C probe, the exception
baseline and the in-situ instrument (Sections~\ref{sec:c-method},
\ref{sec:exc-method} and~\ref{sec:insitu-method}), and the sampling frames of
Section~\ref{sec:frames}.

\subsection{Unit of measurement and interval estimation}

An interval computed over the calls inside one process describes that process's
steady state and is far too narrow: it cannot observe the compilation decisions,
inline-cache states and heap layouts that differ between runs of the same
code~\citep{georges2007statistically}. Conditions in the decomposition of
Section~\ref{sec:rq1} were run in \HostInvocations{} separate processes of
\HostCallsPerInvocation{} calls; conditions in the runtime matrix of
Section~\ref{sec:rq2} in \MatrixInvocations{} processes of
\MatrixCallsPerInvocation{} calls. The largest between-process coefficient of
variation is \HostCvMax\% on the host; in the matrix it is \MatrixCvMax\%,
occurring on the \texttt{\MatrixCvMaxCondition} condition under
\texttt{\MatrixCvMaxEnvironment} rather than across the matrix generally.

Bootstrap intervals over \HostInvocations{} or \MatrixInvocations{} units are
coarse, and we report them as indicative rather than as calibrated coverage.
Where an interval collapses to a point (the timer overhead row of
Table~\ref{tab:host}), it should be read as the clock's resolution showing
through, not as a precise estimate. The per-process values underlying every
interval are in the replication package.

Conditions are interleaved within each round rather than run in blocks. Under
block ordering a thermal excursion or a background process lands entirely on
whichever condition is executing and is indistinguishable from an effect;
interleaved, it is distributed across conditions and appears as between-round
variance. The clock's own cost is measured in the same process as the condition
it accompanies and is reported rather than subtracted: \HostTimerOverhead~$\mu$s
against the smallest quantity reported, \HostCreateSecretKey~$\mu$s.

Differences between conditions are reported throughout as the median of
per-round \emph{paired} differences, not as the difference of two medians. The
two are not identical and the distinction matters at the magnitudes involved;
Section~\ref{sec:residual} and Section~\ref{sec:fix} each rely on it.

\subsection{Taking the JavaScript engine out}

Runtime pairs can separate engine from library only as far as stock builds
allow, so the attribution was tested by removing the engine altogether. We read
the Node source at the two release lines measured
(\texttt{src/crypto/crypto\_keys.cc} at v18.20.8; \texttt{deps/ncrypto} at
v22.23.3, which issues the same sequence) and reproduced in C, call for call,
what \texttt{createPublicKey()} asks of OpenSSL when handed an HMAC secret:
three \texttt{PEM\_bytes\_read\_bio} scans for a public key, an RSA public key
and a certificate, each under an error-queue mark and each followed by
\texttt{BIO\_reset}; then, since none was recognised,
\texttt{PEM\_read\_bio\_PrivateKey} with Node's passphrase callback; then the
C side of Node's error path, which drains the OpenSSL error queue into
strings. The V8 side of that error path is deliberately absent; it is what
Section~\ref{sec:exc-method} measures.

The program was compiled once per OpenSSL release with one compiler
(\CProbeCompiler) and one set of flags, against \CProbeVersionCount{} releases
from \CProbeVersionFirst{} to \CProbeVersionLast{}, each built from source with
identical configuration and linked by run path so that no system library could
be picked up. Each binary reports the version it actually loaded; rows whose
loaded version differed from the one requested would be discarded, and there
were \CProbeIdentityMismatches. The protocol is that of
Section~\ref{sec:unit}: one condition per process, \CProbeInvocations{}
processes per release and condition, \CProbeCallsPerInvocation{} timed calls
each, timed with \texttt{clock\_gettime}, conditions and releases interleaved
within rounds. Besides the full sequence we time its parts separately (the
three public-key scans, the private-key parse, the error-string drain) and a
\emph{successful} parse of an RSA-2048 public key through the same path.

\subsection{Separating the exception from the parse}

Inside Node the failed parse is an OpenSSL failure followed by a V8 exception:
Node constructs an \texttt{Error}, attaches \texttt{library}, \texttt{reason}
and \texttt{code} properties, throws it through the \texttt{createPublicKey()}
frame, and the caller discards it. To measure that exception without the
parse, a JavaScript benchmark builds the same object (message, property
values and error code are captured from a real failure on the runtime being
measured, not typed) and throws it from \ExcDepth{} frames below the catch.
Variants separate its parts: a plain \texttt{Error} thrown and caught in one
frame; the Node-shaped error at depth; the same with stack capture disabled;
and a Node-internal argument error raised before any OpenSSL call. The real
probe is timed alongside, with and without stack capture, so that the stack's
share of it is obtained by paired differencing. All conditions run under
official Node release binaries, which bundle their own OpenSSL, on six
runtimes (\RtEighteenLiteral, \RtTwentyLiteral, \RtTwentyTwoLiteral,
\RtTwentyThreeLiteral, \RtTwentyFourLiteral{} and \RtTwentySixLiteral{}) at
\ExcInvocations{} processes of \ExcCallsPerInvocation{} calls per condition.

Combining this with Section~\ref{sec:c-method} gives a three-way decomposition
of the probe on each runtime: the C cost of the OpenSSL release that runtime
bundles, the V8 exception, and a remainder we attribute to the binding between
them: argument preparation, key-handle allocation and the C++ side of error
construction. The three come from different processes and cannot be paired, so
the interval on the remainder is obtained by resampling each run's processes
independently.

\subsection{In-situ measurement}

The containerised service (Section~\ref{sec:testbed}) exposes a histogram timing the verification call alone. The
instrument is a \emph{paired-scrape difference}: the cumulative histogram is read
before and after each arrival-rate step, and the two readings are differenced,
which isolates that step's distribution from the cumulative counters. Means are
computed as the difference of the histogram's \texttt{\_sum} divided by the
difference of its \texttt{\_count}, so they are exact and do not depend on bucket
boundaries; bucket resolution bounds the reported percentiles only.

Seven arrival rates are driven, and the whole sweep is run twice on one process:
once in ascending order and once in descending order. We report the \emph{range
across the two passes}. This is a range of two, not a confidence interval; no
standard error is computed from it, and it is not written as a symmetric
tolerance. The opposite-order design also serves as a drift test, which
Section~\ref{sec:drift} reports.

One process and two passes is the whole of the in-situ evidence. It is sufficient
to establish that the cost is present in the containerised service and of what order; it is not
sufficient to characterise the shape of the cost-versus-rate relation, and
Section~\ref{sec:insitu} is written accordingly.

\subsection{The three prevalence frames}

\emph{Population co-occurrence counts.} Code-search counts over public
repositories, in four cells: two languages by two import forms. The denominator
of a cell is files importing the library; the numerator is files that \emph{also
mention} the pre-parsing interface. This measures co-occurrence of two strings in
one file, not verified usage. The denominators are returned rounded (each is
an exact multiple of $1024$), so they are magnitudes rather than counts
established file by file.

\emph{A classified sample of call sites.} Call sites located by search,
deduplicated to one file per repository so that no project dominates, with the
key argument traced within its file. This is a convenience sample from a
relevance-ranked, undocumented ranking; inclusion probabilities are unknown and
no weighting back to the population is possible.

\emph{A targeted counter-search.} Every file co-occurring the library with the
pre-parsing interface, fetched in full and adjudicated by hand. This frame is
deliberately biased toward finding the pattern (that is its purpose as a
negative control) and it is therefore disqualified as a prevalence estimator.

\subsection{Host state is recorded per measurement}\label{sec:hoststate}

A service measured on a host it shares with its own load generator, ingress and
sidecars is measured together with whatever else that host is doing. We observed
this directly while building the instrument: re-running one configuration with a
background indexing process consuming host CPU moved the point at which
throughput collapsed by hundreds of requests per second, and both runs are
committed in the replication package~\citep{artifact2026}. A curve that moves
with unrelated host activity does not characterise the service, and this paper
therefore makes no capacity claim. Every in-situ window instead carries the host
load average observed during it, committed beside the measurement, and
Section~\ref{sec:load} uses it.


%% ======================================================================
\section{Additional analyses for RQ1}

This section expands the checks on the decomposition summarised in
Section~\ref{sec:residual}.


\subsection{The parts do not account for all of it}

If the penalty for supplying a string were exactly the discarded attempt plus the
conversion, the two would sum to it. They do not, quite:

\begin{center}
\begin{tabular}{@{}lr@{}}
\toprule
observed penalty              & \DecompObserved~$\mu$s \\
discarded parse $+$ conversion & \DecompPredicted~$\mu$s \\
residual                      & \DecompResidual~$\mu$s \\
\bottomrule
\end{tabular}
\end{center}

All three quantities are medians of per-round paired differences, per
Section~\ref{sec:unit}. The residual's interval,
[\DecompResidualCiLo, \DecompResidualCiHi], excludes zero. The two measured
parts therefore account for \DecompExplainedPct\% and not more; the remainder
lies in the library's own type and claim checking along the non-pre-parsed path.
We report it rather than absorbing it into the headline, because a decomposition
whose parts are declared to sum exactly to the whole has usually not been
checked.

The pre-parsed path carries a smaller unexplained remainder of the same kind, and
we record it for consistency rather than leaving the asymmetry unremarked:
\texttt{verify()} with a pre-parsed key costs \HostJwtPreparsed~$\mu$s against
\HostDecodeOnly~$\mu$s of structural decode plus \HostHmacKeyObject~$\mu$s of
HMAC, leaving roughly \PreparsedUnattributed~$\mu$s unattributed. Neither
residual was measured directly; both are the subject of conditions we did not
run, and both are stated as unattributed rather than assigned.

\subsection{Confirmation without differencing}

The decomposition reaches its conclusion by subtracting conditions, the same
inference style this paper criticises elsewhere. It is therefore checked against
V8's sampling profiler, which attributes self time to named frames with no
arithmetic on our part. At \ProfIterations{} iterations per run and
\ProfRepeats{} runs per environment, matched so the percentages are comparable,
the asymmetric key parse accounts for a median \ProfHostPct\% of self time on
the host (range \ProfHostPctLo--\ProfHostPctHi{} across runs),
\ProfEighteenPct\% (\ProfEighteenPctLo--\ProfEighteenPctHi) under Node
\NodeEighteenNodeMajor, and \ProfTwentySixPct\%
(\ProfTwentySixPctLo--\ProfTwentySixPctHi) under Node
\NodeTwentySixMuslNodeMajor. The frame does not appear among the top frames of
the pre-parsed path in any environment.

Two methods, one answer. We rely on the ordering and the dominance rather than on
the exact percentages: sampling attribution depends on rate and on what else the
process is doing.


%% ======================================================================
\section{Additional analyses for RQ2}

This section expands Sections~\ref{sec:cannot-separate} and~\ref{sec:rq2-arch}.


\subsection{Engine major version does not distinguish the two slow runtimes}

Node \NodeEighteenNodeMajor{} and Node \NodeTwentyNodeMajor{} ship V8
\NodeEighteenVeight{} and \NodeTwentyVeight{} (different major versions of the
engine) and share the OpenSSL \NodeEighteenOpensslSeries{} series. Both pay of
order \NodeEighteenProbeThrows--\NodeTwentyProbeThrows~$\mu$s. Node
\NodeTwentySixMuslNodeMajor{} ships OpenSSL \NodeTwentySixMuslOpensslSeries{} and
pays \NodeTwentySixMuslProbeThrows~$\mu$s on the same libc. Within the OpenSSL
\NodeEighteenOpensslSeries{} series, then, changing the engine's major version
does not move the quantity; had Node \NodeTwentyNodeMajor{} been fast, the
attribution would have been to the engine.

\subsection{What this design cannot separate}

That argument establishes a negative within one OpenSSL series. It does not
establish the positive attribution for the change between series, and we state
the limitation rather than leaving it implicit.

The fast rows are Node \NodeTwentySixMuslNodeMajor. They carry OpenSSL
\NodeTwentySixMuslOpensslSeries{} \emph{and} V8 \NodeTwentySixMuslVeight{}
\emph{and} every other change in that build. OpenSSL version is therefore
perfectly confounded with engine version and with build configuration across the
\NodeEighteenOpensslSeries--\NodeTwentySixMuslOpensslSeries{} boundary. Nothing
in this matrix can separate them. The absence of variation between V8
\NodeEighteenVeight{} and V8 \NodeTwentyVeight{} does not license the inference
that V8 \NodeTwentySixMuslVeight{} is also irrelevant; that is a different
comparison, and we did not make it.

What the matrix supports on its own is the following, and no more: the
magnitude co-varies with the bundled OpenSSL version across four images; the two
slow images differ in engine major version and do not differ in magnitude; and
the two libc variants of one image differ by less than the difference between
series. Separating OpenSSL from the engine requires a crossed design: a
runtime pairing a newer engine with an older OpenSSL, or the parse measured
outside Node entirely. Section~\ref{sec:rq2-causal} reports both, and it is on
them, not on this matrix, that the causal wording used in this paper rests.

\subsection{The effect is confined to key parsing}

Across the four container rows, pre-parsed validation spans
\SpanJwtPreparsed$\times$, the HMAC primitive \SpanHmacString$\times$ and
structural decode \SpanDecodeOnly$\times$. The discarded parse spans
\SpanProbeThrows$\times$. The successful parse moves with it
(\SpanProbeSucceeds$\times$), so this is the cost of key parsing in general
rather than anything peculiar to the failure path. The failure is simply the case
in which an application pays that cost on every request for a result it discards.

\subsection{What this design does not isolate}

A reader may notice that the uncontainerised host row is of the same order as the
Node \NodeTwentySixGlibcNodeMajor{} container row. That comparison is confounded
and we draw no conclusion about containerisation from it: the host runs a
different operating system from the containers, so operating system, C library
and Node build all change alongside the container boundary. Isolating
containerisation would require a Linux-native baseline, which we did not obtain.
The claims above concern only the four container rows, among which
containerisation is constant.

\subsection{A second architecture}

Table~\ref{tab:arch} repeats the runtime matrix, with the same images and
scripts, on x86\_64 (Section~\ref{sec:arch-method}). Absolute values are not
compared across the two machines; the question is whether the ordering holds.

\begin{table*}[!t]
\caption{The runtime matrix of Table~\ref{tab:matrix} on two architectures,
same images and scripts, in $\mu$s. ARM64: \MatrixInvocations{} processes of
\MatrixCallsPerInvocation{} calls. x86\_64, a shared cloud virtual machine:
\XmInvocations{} processes of \XmCallsPerInvocation{} calls. Absolute values
are not comparable across the two machines.}
\label{tab:arch}
\small\setlength{\tabcolsep}{4pt}
\begin{tabular}{@{}llrrrr@{}}
\toprule
 & & \multicolumn{2}{c}{Discarded parse} & \multicolumn{2}{c}{Pre-parsed \texttt{verify()}} \\
Image & OpenSSL & ARM64 & x86\_64 & ARM64 & x86\_64 \\
\midrule
\texttt{node:18.20.8-alpine} & \NodeEighteenOpensslLiteral & \NodeEighteenProbeThrows & \XmEighteenProbeThrows & \NodeEighteenJwtPreparsed & \XmEighteenJwtPreparsed \\
\texttt{node:20-alpine} & \NodeTwentyOpensslLiteral & \NodeTwentyProbeThrows & \XmTwentyProbeThrows & \NodeTwentyJwtPreparsed & \XmTwentyJwtPreparsed \\
\texttt{node:26.6.0-alpine} & \NodeTwentySixMuslOpensslLiteral & \NodeTwentySixMuslProbeThrows & \XmTwentySixMuslProbeThrows & \NodeTwentySixMuslJwtPreparsed & \XmTwentySixMuslJwtPreparsed \\
\texttt{node:26.6.0-bookworm} & \NodeTwentySixGlibcOpensslLiteral & \NodeTwentySixGlibcProbeThrows & \XmTwentySixGlibcProbeThrows & \NodeTwentySixGlibcJwtPreparsed & \XmTwentySixGlibcJwtPreparsed \\
\bottomrule
\end{tabular}
\end{table*}

The property the attribution rests on (every OpenSSL~3.0 image slower than
every OpenSSL~3.5 image) \XmSeriesSplitHolds{} on both architectures. The
strict rank order does not: the two OpenSSL~3.0 images, Node
\NodeEighteenNodeMajor{} and Node \NodeTwentyNodeMajor, exchange places between
ARM64 and x86\_64. Their difference is small against the gap between series on
either machine, and no claim in this paper depends on their relative order. On x86\_64 the discarded parse spans
\XmSpanProbeThrows$\times$ across the images while pre-parsed validation spans
\XmSpanJwtPreparsed$\times$, against \SpanProbeThrows$\times$ and
\SpanJwtPreparsed$\times$ on ARM64. The host decomposition of
Table~\ref{tab:host} was repeated too (Node \XhNode, OpenSSL \XhOpenssl): the
discarded parse is \XhProbeThrows~$\mu$s of a \XhJwtString~$\mu$s call
(\XhProbeShareOfCall\%), the symmetric constructor \XhCreateSecretKey~$\mu$s
and the HMAC primitive \XhHmacString~$\mu$s, so the component ordering of
Table~\ref{tab:host} holds there as well. The full x86\_64 host table is in the
replication package.


%% ======================================================================
\section{Additional analyses for RQ3}

This section tabulates the in-situ sweep plotted in Figure~\ref{fig:insitu},
analyses the difference between its two passes and the role of host load, and
gives the full prevalence counts summarised in Section~\ref{sec:prevalence}.


\begin{table*}[!htbp]
\caption{In-situ validation cost with a string secret, by arrival rate. The two
columns are the ascending and the descending pass over the same seven rates on
one process; each is a single window, and the pair is two observations rather
than an interval. Host load average during each window is committed with the
measurement.}
\label{tab:insitu}
\begin{tabular*}{\textwidth}{@{\extracolsep{\fill}}lrr@{}}
\toprule
Arrival rate & \multicolumn{2}{c}{Mean per call ($\mu$s)} \\
(requests/s) & ascending & descending \\
\midrule
\InsituRateA & \InsituStringAscA & \InsituStringDescA \\
\InsituRateB & \InsituStringAscB & \InsituStringDescB \\
\InsituRateC & \InsituStringAscC & \InsituStringDescC \\
\InsituRateD & \InsituStringAscD & \InsituStringDescD \\
\InsituRateE & \InsituStringAscE & \InsituStringDescE \\
\InsituRateF & \InsituStringAscF & \InsituStringDescF \\
\InsituRateG & \InsituStringAscG & \InsituStringDescG \\
\bottomrule
\end{tabular*}
\end{table*}

\paragraph{The decline with arrival rate is not interpreted} Cost falls
monotonically from \InsituStringAscA~$\mu$s at \InsituRateA~requests/s to
\InsituStringAscG~$\mu$s at \InsituRateG~requests/s. This is the opposite of what
contention alone would produce, and several mechanisms unrelated to validation
would produce it: a fixed per-window overhead amortising over more calls,
processor frequency scaling and idle-state exit latency at low rates, cache and
branch-predictor residency, or continued tier-up in the engine. We did not run
the conditions that would distinguish them. We therefore report the numbers and
do not claim that validation cost depends on arrival rate; establishing that
would require varying window duration at fixed rate, pinning the frequency
governor, and extending warmup at low rates, none of which we did.

\subsection{The two passes differ systematically}\label{sec:drift}

Sweeping up and then down is not only replication; it detects warm-up, thermal
and hysteresis effects. If the difference between passes were noise, its sign
would alternate across rate points. It does not. In the string-secret sweep the
descending pass is slower at \InsituStringSameSign{} of \InsituStringPoints{}
rate points; in the pre-parsed sweep the descending pass is \emph{faster} at
\InsituPreparsedSameSign{} of \InsituPreparsedPoints. A consistent sign across
seven points is not what noise produces.

This has a direct consequence for how the ranges in Table~\ref{tab:insitu} should
be read. A systematic component means the two passes are not independent samples
of a common value, so the range between them understates uncertainty further than
a range of two already does. We therefore report ranges and draw no inference that
depends on their width. It is also a methodological finding in its own right: a
single-direction sweep would have produced a curve carrying a consistent bias with
no way to detect it.

\subsection{Host load, not the condition, drives the outliers}\label{sec:load}

Agreement between passes varies considerably across windows. We distinguish
\emph{relative} disagreement, the gap between passes as a fraction of their mean,
from \emph{absolute} disagreement, the gap in microseconds; the two behave
differently and the distinction carries this subsection.

The widest relative disagreement in the dataset and the tightest occur at the same
arrival rate in different sweeps. The wide one is the highest-load window in its
own sweep on the descending pass, though not on the ascending pass, whose maximum
load falls at a different rate. The tight one, at the same arrival rate in another
sweep, sits at a lower load on both passes and its two passes agree to within
\InsituTightestAgreement~$\mu$s. Relative disagreement therefore does not track
the condition under measurement.

We considered two alternative explanations for the relative asymmetry and reject
both. A small-mean effect would require absolute disagreement to be roughly
constant, so that its relative size grew only as means shrank; absolute
disagreement instead varies over three orders of magnitude and is largest where
the means are largest, which is the opposite pattern. A measurement floor would
require the affected values to sit at the instrument's resolution; they do not,
and in any case the means reported here are computed from the histogram's sum and
count rather than from its buckets, so bucket resolution cannot bound them.

The operative consequence is that in-situ figures are reported with the host load
recorded during their window, and comparisons between windows at different loads
are not treated as comparisons between conditions.

\paragraph{Population co-occurrence} Of \PopJsRequireFiles{} public JavaScript
files importing the library through \texttt{require},
\PopJsRequireFilesWithApi{} also mention the pre-parsing interface; for
JavaScript through \texttt{import} the figures are \PopJsImportFiles{} and
\PopJsImportFilesWithApi; for TypeScript through \texttt{require},
\PopTsRequireFiles{} and \PopTsRequireFilesWithApi; and for TypeScript through
\texttt{import}, \PopTsImportFiles{} and \PopTsImportFilesWithApi. Cells are not
pooled into a rate: a file using both import forms would be counted twice, so
their sum \PopSumFiles{} is an upper bound on the union and the largest cell
\PopLargestCellFiles{} a lower bound. No cell exceeds \PopWorstCellPct\% of files
mentioning the interface. These are search-engine estimates of file counts over
public repositories only, they measure co-occurrence rather than use, and the
denominators are returned \emph{rounded} (each is an exact multiple of $1024$), so they should be read as magnitudes rather than as counts established file
by file.

Matching only the conventional \texttt{jwt.verify()} spelling would filter the
corpus toward one import idiom, so call sites are found by resolving the local
binding of the import, covering arbitrary aliases, namespace imports,
destructured and renamed imports, and direct calls on the require expression.
That binding resolution contributes \SampleRematchAdded{} of the call sites
above that the conventional spelling misses, of which \RematchSitesPreparsed{}
pass a pre-parsed key.


%% ======================================================================
\section{Extended threats to validity}

This section gives the full discussion that Section~\ref{sec:threats}
summarises.


The limitations below are grouped by what they bear on: the measurements,
the deployment evidence, the survey of public code, and the scope of the
claims. Each states what the paper does not claim as a result.

\subsection{Measurement}\label{sec:threats-measure}

\paragraph{Two architectures, one of them shared} The primary measurements
come from a single ARM64 host. Section~\ref{sec:rq2-arch} repeats them on
x86\_64, and there the split between OpenSSL series and the component
ordering of the host decomposition hold. But the x86\_64 machine is a virtual machine on shared
cloud infrastructure, whose absolute timings are known to be
unreliable~\citep{laaber2019cloud}; we therefore take only orderings and ratios
between interleaved conditions from it, and no magnitude in this paper should
be read as holding on x86\_64 hardware in general. A repetition on a dedicated
x86\_64 host, which the replication package scripts in one command, would turn
the ordering check into a replication of magnitudes. The C probe, the exception
baseline and the cross-library comparison were measured only on the x86\_64
machine. Their conclusions rest on ratios between conditions measured in the
same rounds, which is the kind of quantity a shared machine can support, and
the effects they report are far larger than the machine's noise: the median
between-process coefficient of variation is \CProbeCvMedian\%,
\ExcCvMedian\% and \XlibCvMedian\% in the three runs. The largest values
(\CProbeCvMax\%, \ExcCvMax\% and \XlibCvMax\%) occur on sub-microsecond or
JIT-dominated conditions that carry no claim, and medians of per-process
medians are insensitive to them. The C probe needs no JavaScript runtime and
the scripts rebuild every OpenSSL release from source, so repeating it on the
ARM64 host or a dedicated x86\_64 machine is one command.

\paragraph{The C probe reproduces Node's calls, not Node's build}
Section~\ref{sec:rq2-causal} issues the OpenSSL calls Node issues, as read from
the Node source, against OpenSSL releases we built. Node builds its bundled
OpenSSL with its own configuration, so the C figures approximate rather than
reproduce the parse cost inside a given runtime, and the binding column of
Table~\ref{tab:decomp} inherits that approximation; it is slightly negative
on two OpenSSL~3.0 runtimes for that reason. The causal claim does not inherit
it: it compares releases under one build configuration with nothing else
varying.

\paragraph{The exception baseline is a reconstruction} Node constructs the
error from C++ through V8's API; our baseline constructs an error of the same
shape, message and properties from JavaScript at a comparable stack depth. The
two paths share V8's error construction and stack capture, and the paired
stack-capture measurement on the real probe agrees in order of magnitude with
the synthetic one, but they are not the same code path. The ARM64 host figure
of Table~\ref{tab:host} was not itself decomposed; Table~\ref{tab:decomp}
decomposes the same component on six other runtimes.

\paragraph{The headline host figure has no per-row OpenSSL version} The host
decomposition of Section~\ref{sec:rq1} predates the per-measurement OpenSSL
field. The version in force is known from the run timestamp and a note recorded
at the time: OpenSSL \HostOpensslInferred, which a package upgrade later that
day moved to \HostOpensslAfterDrift{} under the same Node binary. We report it
as inferred, marked as such in Table~\ref{tab:matrix}, rather than as recorded.
Three measurements that do carry their version bound how much this matters.
The same probe costs \RtTwentySixProbe~$\mu$s under Node \RtTwentySixLiteral{}
with OpenSSL \RtTwentySixOpenssl, \XhProbeThrows~$\mu$s in the x86\_64 host
repetition with OpenSSL \XhOpenssl, and \NodeTwentySixGlibcProbeThrows~$\mu$s in
the Node \NodeTwentySixGlibcNodeMajor{} container with OpenSSL
\NodeTwentySixGlibcOpensslLiteral{}: all of the same order as
\HostProbeThrows~$\mu$s, and all from OpenSSL releases after the step of
Table~\ref{tab:cprobe}. The headline figure is therefore representative of a
post-3.2 OpenSSL, which is all the paper uses it for; it would be wrong by more
than an order of magnitude only if the host had run an OpenSSL~3.0 or~3.1
release, which the record excludes.

\subsection{Deployment measurement}\label{sec:threats-deploy}

\paragraph{Ranges of two, with a systematic component} The in-situ dispersion is
the range across an ascending and a descending pass. It is not a confidence
interval, and Section~\ref{sec:drift} shows the two passes differ systematically
rather than randomly, so they are not independent samples of a common value and
the range understates uncertainty further. No inference in this paper depends on
the width of those ranges.

\paragraph{One in-situ process, two passes} The rate sweep rests on one
process sweeping seven rates twice. That is sufficient for the claim made (the cost is present in the containerised service and of this order) and insufficient for any
claim about the shape of the cost-versus-rate relation, which
Section~\ref{sec:insitu} accordingly does not make. The magnitude claim does not
rest on the sweep alone: the A/B of Section~\ref{sec:abcorroboration} is a
separate pair of windows, measured with a separate instrument, and it agrees
with the microbenchmark prediction on processor time. The wall-clock excess over
that prediction is attributed to time off the processor and not decomposed
further; the decline with arrival rate has several candidate explanations
unrelated to validation that we did not test (Section~\ref{sec:insitu}).

\paragraph{Host load, and the confound that motivated recording it} A
throughput curve measured on a host shared with its own load generator moves with
unrelated activity on that host, as Section~\ref{sec:hoststate} describes, which
is why no capacity claim appears in this paper. Every in-situ window here carries
the load average observed during it, and Section~\ref{sec:load} shows that window-to-window
agreement tracks that load rather than the condition. Contention inflates a
duration rather than deflating it, so in-situ figures are upper bounds; but
comparisons between windows at different loads are not comparisons between
conditions, and we do not treat them as such.

\paragraph{The A/B is two windows at unequal load} The two arms of
Table~\ref{tab:ab} are consecutive windows rather than interleaved rounds, and
their host load is known only at the instant each began. That bars comparing
them on wall-clock time, and we do not. It does not bar the CPU comparison the
section rests on: contention inflates elapsed time, whereas a container's own
cgroup accounting charges it only the processor time it used. Attaching the
per-second host sampler to a re-run of both arms, interleaved, would remove the
limitation; we did not.

\subsection{Survey of public code}\label{sec:threats-survey}

\paragraph{Sampling frames are not interchangeable} The three prevalence
instruments of Section~\ref{sec:frames} have different units and different
selection mechanisms, and none supports a population proportion. The
co-occurrence counts measure two strings appearing in one file, with a
false-positive rate established at \AdjudicatedFalsePositive{} of
\AdjudicatedTotal{} on hand inspection. The classified sample has unknown
inclusion probabilities, a limitation of repository search noted more generally
by \citet{kalliamvakou2014promises}. The counter-search is
selected on the outcome.

\paragraph{The matcher's coverage} Call sites are found by resolving the local
binding of the import, which covers aliases, namespace imports, destructured and
renamed imports and direct calls on the require expression. It does not resolve
bindings across files, so a project that constructs a pre-parsed key in one file
and validates in another is invisible to this analysis. This is the strongest
threat to Section~\ref{sec:prevalence}; within-file tracing found no such case,
which is weak evidence against it.

\paragraph{One rater} The \AdjudicatedTotal{} automatically flagged call sites
were read and classified by the first author, and the authors are not
disinterested: the paper's claim is stronger the more of them are false
positives, and \AdjudicatedFalsePositive{} were classified that way. There was no
second rater and no inter-rater agreement to report. The per-case verdicts, the
evidence for each and the hashes of the files they were read from are published
so the adjudication can be repeated, which is weaker than having repeated it.

\subsection{Scope}\label{sec:threats-scope}

\paragraph{Token caching is not in the comparison} A service that caches
verified tokens pays this cost once per token rather than once per request. We
did not measure such a configuration, so the per-request figures here do not
describe a caching deployment.

\paragraph{The security analysis is from source} Section~\ref{sec:security}
establishes reachability and provenance by reading \texttt{verify.js} at the
releases named and by running the library's own tests; it does not drive
forged tokens at the containerised service, time the patched library on tokens
declaring a non-HMAC algorithm, or time version 8.5.1. Its statements are
therefore about which inputs reach the attempt and when the code arrived, not
about throughput under attack, and the paper makes no availability claim.

\paragraph{Seven libraries, one version each} Section~\ref{sec:rq4} covers
seven libraries in four languages, chosen as the most used in each ecosystem,
each at one release. Other libraries, and other releases of these, may resolve
keys differently; each classification is from source at the version measured.
The comparison supports the claim that the cost follows a design choice rather
than JWT validation; it does not estimate how many libraries make that
choice.

%% ======================================================================
\section{Relocation to the service mesh: a design note}


\subsection{Relocation is not the first remedy}\label{sec:offload}

A cost believed to be cryptographic invites relocation to a sidecar or gateway.
The cost measured here is not cryptographic and is largely avoidable in place,
so an architectural remedy that also moves a trust boundary should not be the
first thing reached for. We implemented that relocation for the service
studied here; applying a mesh validation policy duplicates verification rather
than moving it, and making it effective requires the application to trust a
proxy-supplied claims header, a known anti-pattern that needs a filter
stripping client-supplied copies. We report that only as a design note: it
carries no measurement, and the configuration, the observation and the
inconclusive automated probe are documented in the replication
package~\citep{artifact2026}.



\bibliographystyle{elsarticle-harv}
\bibliography{references}
\end{document}
